% Options for packages loaded elsewhere
\PassOptionsToPackage{unicode}{hyperref}
\PassOptionsToPackage{hyphens}{url}
%
\documentclass[
]{article}
\usepackage{amsmath,amssymb}
\usepackage{iftex}
\ifPDFTeX
  \usepackage[T1]{fontenc}
  \usepackage[utf8]{inputenc}
  \usepackage{textcomp} % provide euro and other symbols
\else % if luatex or xetex
  \usepackage{unicode-math} % this also loads fontspec
  \defaultfontfeatures{Scale=MatchLowercase}
  \defaultfontfeatures[\rmfamily]{Ligatures=TeX,Scale=1}
\fi
\usepackage{lmodern}
\ifPDFTeX\else
  % xetex/luatex font selection
\fi
% Use upquote if available, for straight quotes in verbatim environments
\IfFileExists{upquote.sty}{\usepackage{upquote}}{}
\IfFileExists{microtype.sty}{% use microtype if available
  \usepackage[]{microtype}
  \UseMicrotypeSet[protrusion]{basicmath} % disable protrusion for tt fonts
}{}
\makeatletter
\@ifundefined{KOMAClassName}{% if non-KOMA class
  \IfFileExists{parskip.sty}{%
    \usepackage{parskip}
  }{% else
    \setlength{\parindent}{0pt}
    \setlength{\parskip}{6pt plus 2pt minus 1pt}}
}{% if KOMA class
  \KOMAoptions{parskip=half}}
\makeatother
\usepackage{xcolor}
\usepackage[margin=1in]{geometry}
\usepackage{color}
\usepackage{fancyvrb}
\newcommand{\VerbBar}{|}
\newcommand{\VERB}{\Verb[commandchars=\\\{\}]}
\DefineVerbatimEnvironment{Highlighting}{Verbatim}{commandchars=\\\{\}}
% Add ',fontsize=\small' for more characters per line
\usepackage{framed}
\definecolor{shadecolor}{RGB}{248,248,248}
\newenvironment{Shaded}{\begin{snugshade}}{\end{snugshade}}
\newcommand{\AlertTok}[1]{\textcolor[rgb]{0.94,0.16,0.16}{#1}}
\newcommand{\AnnotationTok}[1]{\textcolor[rgb]{0.56,0.35,0.01}{\textbf{\textit{#1}}}}
\newcommand{\AttributeTok}[1]{\textcolor[rgb]{0.13,0.29,0.53}{#1}}
\newcommand{\BaseNTok}[1]{\textcolor[rgb]{0.00,0.00,0.81}{#1}}
\newcommand{\BuiltInTok}[1]{#1}
\newcommand{\CharTok}[1]{\textcolor[rgb]{0.31,0.60,0.02}{#1}}
\newcommand{\CommentTok}[1]{\textcolor[rgb]{0.56,0.35,0.01}{\textit{#1}}}
\newcommand{\CommentVarTok}[1]{\textcolor[rgb]{0.56,0.35,0.01}{\textbf{\textit{#1}}}}
\newcommand{\ConstantTok}[1]{\textcolor[rgb]{0.56,0.35,0.01}{#1}}
\newcommand{\ControlFlowTok}[1]{\textcolor[rgb]{0.13,0.29,0.53}{\textbf{#1}}}
\newcommand{\DataTypeTok}[1]{\textcolor[rgb]{0.13,0.29,0.53}{#1}}
\newcommand{\DecValTok}[1]{\textcolor[rgb]{0.00,0.00,0.81}{#1}}
\newcommand{\DocumentationTok}[1]{\textcolor[rgb]{0.56,0.35,0.01}{\textbf{\textit{#1}}}}
\newcommand{\ErrorTok}[1]{\textcolor[rgb]{0.64,0.00,0.00}{\textbf{#1}}}
\newcommand{\ExtensionTok}[1]{#1}
\newcommand{\FloatTok}[1]{\textcolor[rgb]{0.00,0.00,0.81}{#1}}
\newcommand{\FunctionTok}[1]{\textcolor[rgb]{0.13,0.29,0.53}{\textbf{#1}}}
\newcommand{\ImportTok}[1]{#1}
\newcommand{\InformationTok}[1]{\textcolor[rgb]{0.56,0.35,0.01}{\textbf{\textit{#1}}}}
\newcommand{\KeywordTok}[1]{\textcolor[rgb]{0.13,0.29,0.53}{\textbf{#1}}}
\newcommand{\NormalTok}[1]{#1}
\newcommand{\OperatorTok}[1]{\textcolor[rgb]{0.81,0.36,0.00}{\textbf{#1}}}
\newcommand{\OtherTok}[1]{\textcolor[rgb]{0.56,0.35,0.01}{#1}}
\newcommand{\PreprocessorTok}[1]{\textcolor[rgb]{0.56,0.35,0.01}{\textit{#1}}}
\newcommand{\RegionMarkerTok}[1]{#1}
\newcommand{\SpecialCharTok}[1]{\textcolor[rgb]{0.81,0.36,0.00}{\textbf{#1}}}
\newcommand{\SpecialStringTok}[1]{\textcolor[rgb]{0.31,0.60,0.02}{#1}}
\newcommand{\StringTok}[1]{\textcolor[rgb]{0.31,0.60,0.02}{#1}}
\newcommand{\VariableTok}[1]{\textcolor[rgb]{0.00,0.00,0.00}{#1}}
\newcommand{\VerbatimStringTok}[1]{\textcolor[rgb]{0.31,0.60,0.02}{#1}}
\newcommand{\WarningTok}[1]{\textcolor[rgb]{0.56,0.35,0.01}{\textbf{\textit{#1}}}}
\usepackage{graphicx}
\makeatletter
\def\maxwidth{\ifdim\Gin@nat@width>\linewidth\linewidth\else\Gin@nat@width\fi}
\def\maxheight{\ifdim\Gin@nat@height>\textheight\textheight\else\Gin@nat@height\fi}
\makeatother
% Scale images if necessary, so that they will not overflow the page
% margins by default, and it is still possible to overwrite the defaults
% using explicit options in \includegraphics[width, height, ...]{}
\setkeys{Gin}{width=\maxwidth,height=\maxheight,keepaspectratio}
% Set default figure placement to htbp
\makeatletter
\def\fps@figure{htbp}
\makeatother
\setlength{\emergencystretch}{3em} % prevent overfull lines
\providecommand{\tightlist}{%
  \setlength{\itemsep}{0pt}\setlength{\parskip}{0pt}}
\setcounter{secnumdepth}{-\maxdimen} % remove section numbering
\ifLuaTeX
  \usepackage{selnolig}  % disable illegal ligatures
\fi
\usepackage{bookmark}
\IfFileExists{xurl.sty}{\usepackage{xurl}}{} % add URL line breaks if available
\urlstyle{same}
\hypersetup{
  pdftitle={network\_analysis},
  pdfauthor={Pietro Olivetti},
  hidelinks,
  pdfcreator={LaTeX via pandoc}}

\title{network\_analysis}
\author{Pietro Olivetti}
\date{2024-10-17}

\begin{document}
\maketitle

\section{Network analysis of Depression and Anxiety
Symptomatology}\label{network-analysis-of-depression-and-anxiety-symptomatology}

\subsection{Packages and loading}\label{packages-and-loading}

\begin{Shaded}
\begin{Highlighting}[]
\FunctionTok{library}\NormalTok{(dplyr)}
\FunctionTok{library}\NormalTok{(igraph)}
\FunctionTok{library}\NormalTok{(qgraph)}
\FunctionTok{library}\NormalTok{(bootnet)}
\FunctionTok{library}\NormalTok{(corpcor)}
\FunctionTok{library}\NormalTok{(openxlsx)}
\FunctionTok{library}\NormalTok{(ggplot2)}
\FunctionTok{library}\NormalTok{(tidyr)}
\FunctionTok{library}\NormalTok{(matrixcalc)}
\FunctionTok{library}\NormalTok{(e1071)}
\FunctionTok{library}\NormalTok{(networktools)}
\FunctionTok{library}\NormalTok{(EGAnet)}
\CommentTok{\#library(smacof)}
\FunctionTok{source}\NormalTok{(}\StringTok{\textquotesingle{}../data/questions.R\textquotesingle{}}\NormalTok{)}
\end{Highlighting}
\end{Shaded}

\begin{Shaded}
\begin{Highlighting}[]
\NormalTok{multiC }\OtherTok{\textless{}{-}} \ControlFlowTok{function}\NormalTok{(cor\_mat, }\AttributeTok{t\_hold =} \FloatTok{0.8}\NormalTok{, }\AttributeTok{name =} \StringTok{""}\NormalTok{)\{threshold }\OtherTok{\textless{}{-}}\NormalTok{ t\_hold}
\NormalTok{high\_corr\_pairs }\OtherTok{\textless{}{-}} \FunctionTok{which}\NormalTok{(}\FunctionTok{abs}\NormalTok{(cor\_mat) }\SpecialCharTok{\textgreater{}}\NormalTok{ threshold }\SpecialCharTok{\&} \FunctionTok{abs}\NormalTok{(cor\_mat) }\SpecialCharTok{\textless{}} \DecValTok{1}\NormalTok{, }\AttributeTok{arr.ind =} \ConstantTok{TRUE}\NormalTok{)}

\CommentTok{\# variable names and corresponding correlation values}
\NormalTok{high\_corr\_list }\OtherTok{\textless{}{-}} \FunctionTok{data.frame}\NormalTok{(}
  \AttributeTok{Variable1 =} \FunctionTok{rownames}\NormalTok{(cor\_mat)[high\_corr\_pairs[, }\DecValTok{1}\NormalTok{]],}
  \AttributeTok{Variable2 =} \FunctionTok{colnames}\NormalTok{(cor\_mat)[high\_corr\_pairs[, }\DecValTok{2}\NormalTok{]],}
  \AttributeTok{Correlation =}\NormalTok{ cor\_mat[high\_corr\_pairs]}
\NormalTok{)}

\CommentTok{\# variables and their correlation values}
\NormalTok{df }\OtherTok{\textless{}{-}}\NormalTok{ high\_corr\_list }\SpecialCharTok{\%\textgreater{}\%}
  \FunctionTok{arrange}\NormalTok{(}\FunctionTok{desc}\NormalTok{(Correlation))\}}
\end{Highlighting}
\end{Shaded}

\subsection{Organizing data}\label{organizing-data}

\subsubsection{All variables}\label{all-variables}

Creating the dataframe and the group's list for all variables

\begin{Shaded}
\begin{Highlighting}[]
\CommentTok{\# Dataframe}
\NormalTok{df }\OtherTok{\textless{}{-}} \FunctionTok{read.csv}\NormalTok{(}\StringTok{\textquotesingle{}../data/data\_cleaned.csv\textquotesingle{}}\NormalTok{)}

\CommentTok{\# Groups}
\NormalTok{groups\_all\_variables }\OtherTok{\textless{}{-}} \FunctionTok{list}\NormalTok{(}\AttributeTok{Ids =} \DecValTok{1}\NormalTok{, }\AttributeTok{Sociodem =} \DecValTok{2}\SpecialCharTok{:}\DecValTok{8}\NormalTok{, }\AttributeTok{Pers\_BFI2 =} \DecValTok{9}\SpecialCharTok{:}\DecValTok{23}\NormalTok{, }\AttributeTok{Pers\_hexaco =} \DecValTok{24}\SpecialCharTok{:}\DecValTok{33}\NormalTok{, }\AttributeTok{Depr\_cesd =} \DecValTok{34}\SpecialCharTok{:}\DecValTok{53}\NormalTok{, }\AttributeTok{Depr\_phq =} \DecValTok{54}\SpecialCharTok{:}\DecValTok{62}\NormalTok{, }\AttributeTok{Depr\_promis =} \DecValTok{63}\SpecialCharTok{:}\DecValTok{90}\NormalTok{, }\AttributeTok{Depr\_qids =} \DecValTok{91}\SpecialCharTok{:}\DecValTok{106}\NormalTok{, }\AttributeTok{Depr\_Ans\_dass =} \DecValTok{107}\SpecialCharTok{:}\DecValTok{120}\NormalTok{, }\AttributeTok{Depr\_Ans\_hads =} \DecValTok{121}\SpecialCharTok{:}\DecValTok{134}\NormalTok{, }\AttributeTok{Ans\_gad =} \DecValTok{135}\SpecialCharTok{:}\DecValTok{141}\NormalTok{, }\AttributeTok{Ans\_tmas =} \DecValTok{142}\SpecialCharTok{:}\DecValTok{191}\NormalTok{, }\AttributeTok{Ans\_sas =} \DecValTok{192}\SpecialCharTok{:}\DecValTok{206}\NormalTok{)}

\DocumentationTok{\#\#\#\#\#\#\#\#}
\end{Highlighting}
\end{Shaded}

\subsubsection{All Scales}\label{all-scales}

Creating the dataframe and the group's list for all scales

\begin{Shaded}
\begin{Highlighting}[]
\CommentTok{\# Dataframe}
\NormalTok{df\_all\_scales }\OtherTok{\textless{}{-}}\NormalTok{ df[}\DecValTok{9}\SpecialCharTok{:}\DecValTok{206}\NormalTok{]}

\NormalTok{groups\_all\_scales }\OtherTok{\textless{}{-}} \FunctionTok{list}\NormalTok{(}\AttributeTok{Pers\_BFI2 =} \DecValTok{1}\SpecialCharTok{:}\DecValTok{15}\NormalTok{, }\AttributeTok{Pers\_hexaco =} \DecValTok{16}\SpecialCharTok{:}\DecValTok{25}\NormalTok{, }\AttributeTok{Depr\_cesd =} \DecValTok{26}\SpecialCharTok{:}\DecValTok{45}\NormalTok{, }\AttributeTok{Depr\_phq =} \DecValTok{46}\SpecialCharTok{:}\DecValTok{54}\NormalTok{, }\AttributeTok{Depr\_promis =} \DecValTok{55}\SpecialCharTok{:}\DecValTok{82}\NormalTok{, }\AttributeTok{Depr\_qids =} \DecValTok{83}\SpecialCharTok{:}\DecValTok{98}\NormalTok{, }\AttributeTok{Depr\_Ans\_dass =} \DecValTok{99}\SpecialCharTok{:}\DecValTok{112}\NormalTok{, }\AttributeTok{Depr\_Ans\_hads =} \DecValTok{113}\SpecialCharTok{:}\DecValTok{126}\NormalTok{, }\AttributeTok{Ans\_gad =} \DecValTok{127}\SpecialCharTok{:}\DecValTok{133}\NormalTok{, }\AttributeTok{Ans\_tmas =} \DecValTok{134}\SpecialCharTok{:}\DecValTok{183}\NormalTok{, }\AttributeTok{Ans\_sas =} \DecValTok{184}\SpecialCharTok{:}\DecValTok{198}\NormalTok{)}

\NormalTok{group\_names }\OtherTok{\textless{}{-}} \FunctionTok{c}\NormalTok{(}\StringTok{"Ids"}\NormalTok{, }\StringTok{"Sociodem"}\NormalTok{, }\StringTok{"Pers\_BFI2"}\NormalTok{, }\StringTok{"Pers\_hexaco"}\NormalTok{, }\StringTok{"Depr\_cesd"}\NormalTok{, }\StringTok{"Depr\_phq"}\NormalTok{, }\StringTok{"Depr\_promis"}\NormalTok{, }\StringTok{"Depr\_qids"}\NormalTok{, }\StringTok{"Depr\_Ans\_dass"}\NormalTok{, }\StringTok{"Depr\_Ans\_hads"}\NormalTok{, }\StringTok{"Ans\_gad"}\NormalTok{, }\StringTok{"Ans\_tmas"}\NormalTok{, }\StringTok{"Ans\_sas"}\NormalTok{)}
\DocumentationTok{\#\#\#\#\#\#\#}
\end{Highlighting}
\end{Shaded}

\subsubsection{Scales related to anxiety and
depression}\label{scales-related-to-anxiety-and-depression}

Creating the dataframe and the group's list for all scales related to
depression and anxiety

\begin{Shaded}
\begin{Highlighting}[]
\CommentTok{\# Dataframe}
\NormalTok{df\_anx\_dp }\OtherTok{\textless{}{-}}\NormalTok{ df[}\DecValTok{34}\SpecialCharTok{:}\DecValTok{206}\NormalTok{]}



\CommentTok{\# Groups}
\NormalTok{groups\_anx\_dep }\OtherTok{\textless{}{-}} \FunctionTok{list}\NormalTok{(}\AttributeTok{D\_cesd =} \DecValTok{1}\SpecialCharTok{:}\DecValTok{20}\NormalTok{, }\AttributeTok{D\_phq =} \DecValTok{21}\SpecialCharTok{:}\DecValTok{29}\NormalTok{, }\AttributeTok{D\_promis =} \DecValTok{30}\SpecialCharTok{:}\DecValTok{57}\NormalTok{, }\AttributeTok{D\_qids =} \DecValTok{58}\SpecialCharTok{:}\DecValTok{73}\NormalTok{, }\AttributeTok{A\_dass =} \DecValTok{74}\SpecialCharTok{:}\DecValTok{80}\NormalTok{ , }\AttributeTok{D\_dass =} \DecValTok{81}\SpecialCharTok{:}\DecValTok{87}\NormalTok{ , }\AttributeTok{A\_hads =} \DecValTok{88}\SpecialCharTok{:}\DecValTok{94}\NormalTok{, }\AttributeTok{D\_hads =} \DecValTok{95}\SpecialCharTok{:}\DecValTok{101}\NormalTok{, }\AttributeTok{A\_gad =} \DecValTok{102}\SpecialCharTok{:}\DecValTok{108}\NormalTok{, }\AttributeTok{A\_sas =} \DecValTok{109}\SpecialCharTok{:}\DecValTok{123}\NormalTok{, }\AttributeTok{A\_tmas =} \DecValTok{124}\SpecialCharTok{:}\DecValTok{173}\NormalTok{)}
\end{Highlighting}
\end{Shaded}

\subsubsection{Scales related to anxiety and depression without
TMAS}\label{scales-related-to-anxiety-and-depression-without-tmas}

TMAS items are excluded because they are binary, unlike the Likert-scale
items in the dataset, which could affect the accuracy of network
analysis by increasing multicollinearity and distorting symptom
relationships. Dichotomizing data, such as converting scores to binary
categories, risks information loss, and while Cai (2022) uses binary
data in network analysis, the potential impact of this on network
structure is not addressed. Studies like Burger et al.~(2022) advise
caution when transforming data for network analysis, particularly with
binarization. Excluding TMAS is a methodological choice that aims to
ensure a more reliable and consistent analysis by using Likert scales
for all symptoms.

\begin{Shaded}
\begin{Highlighting}[]
\CommentTok{\# Dataframe}

\NormalTok{D\_cesd }\OtherTok{\textless{}{-}}\NormalTok{ df\_anx\_dp[}\DecValTok{1}\SpecialCharTok{:}\DecValTok{20}\NormalTok{]}
\NormalTok{D\_phq }\OtherTok{\textless{}{-}}\NormalTok{ df\_anx\_dp[}\DecValTok{21}\SpecialCharTok{:}\DecValTok{29}\NormalTok{]}
\NormalTok{D\_promis }\OtherTok{\textless{}{-}}\NormalTok{ df\_anx\_dp[}\DecValTok{30}\SpecialCharTok{:}\DecValTok{57}\NormalTok{]}
\NormalTok{D\_qids }\OtherTok{\textless{}{-}}\NormalTok{ df\_anx\_dp[}\DecValTok{58}\SpecialCharTok{:}\DecValTok{73}\NormalTok{]}
\NormalTok{A\_dass }\OtherTok{\textless{}{-}}\NormalTok{ df\_anx\_dp[}\FunctionTok{c}\NormalTok{(}\DecValTok{74}\NormalTok{, }\DecValTok{76}\NormalTok{, }\DecValTok{78}\NormalTok{, }\DecValTok{79}\NormalTok{, }\DecValTok{82}\NormalTok{, }\DecValTok{85}\NormalTok{, }\DecValTok{86}\NormalTok{)] }\CommentTok{\#74 ...80}
\NormalTok{D\_dass }\OtherTok{\textless{}{-}}\NormalTok{ df\_anx\_dp[}\FunctionTok{c}\NormalTok{(}\DecValTok{75}\NormalTok{, }\DecValTok{77}\NormalTok{, }\DecValTok{80}\NormalTok{, }\DecValTok{81}\NormalTok{, }\DecValTok{83}\NormalTok{, }\DecValTok{84}\NormalTok{, }\DecValTok{87}\NormalTok{)] }\CommentTok{\#81 ... 87 }
\NormalTok{A\_hads }\OtherTok{\textless{}{-}}\NormalTok{ df\_anx\_dp[}\FunctionTok{c}\NormalTok{(}\DecValTok{88}\NormalTok{,}\DecValTok{90}\NormalTok{,}\DecValTok{92}\NormalTok{,}\DecValTok{94}\NormalTok{,}\DecValTok{96}\NormalTok{, }\DecValTok{98}\NormalTok{, }\DecValTok{100}\NormalTok{)] }\CommentTok{\# 88 ... 94}
\NormalTok{D\_hads }\OtherTok{\textless{}{-}}\NormalTok{ df\_anx\_dp[}\FunctionTok{c}\NormalTok{(}\DecValTok{89}\NormalTok{,}\DecValTok{91}\NormalTok{,}\DecValTok{93}\NormalTok{,}\DecValTok{95}\NormalTok{,}\DecValTok{97}\NormalTok{,}\DecValTok{99}\NormalTok{,}\DecValTok{101}\NormalTok{)] }\CommentTok{\# 95 ... 101}
\NormalTok{A\_gad }\OtherTok{\textless{}{-}}\NormalTok{ df\_anx\_dp[}\DecValTok{102}\SpecialCharTok{:}\DecValTok{108}\NormalTok{]}
\NormalTok{A\_sas }\OtherTok{\textless{}{-}}\NormalTok{ df\_anx\_dp[}\DecValTok{159}\SpecialCharTok{:}\DecValTok{173}\NormalTok{] }\CommentTok{\#109 ... 123}
\NormalTok{A\_tmas }\OtherTok{\textless{}{-}}\NormalTok{ df\_anx\_dp[}\DecValTok{109}\SpecialCharTok{:}\DecValTok{158}\NormalTok{] }\CommentTok{\#124 ... 174}


\NormalTok{df\_anx\_dp\_adj }\OtherTok{\textless{}{-}} \FunctionTok{as.data.frame}\NormalTok{(}\FunctionTok{list}\NormalTok{(D\_cesd, D\_phq, D\_promis, D\_qids, A\_dass, D\_dass, A\_hads, D\_hads, A\_gad, A\_sas))}



\CommentTok{\# Groups}
\NormalTok{groups\_anx\_dep\_adj }\OtherTok{\textless{}{-}} \FunctionTok{list}\NormalTok{(}\AttributeTok{D\_cesd =} \DecValTok{1}\SpecialCharTok{:}\DecValTok{20}\NormalTok{, }\AttributeTok{D\_phq =} \DecValTok{21}\SpecialCharTok{:}\DecValTok{29}\NormalTok{, }\AttributeTok{D\_promis =} \DecValTok{30}\SpecialCharTok{:}\DecValTok{57}\NormalTok{, }\AttributeTok{D\_qids =} \DecValTok{58}\SpecialCharTok{:}\DecValTok{73}\NormalTok{, }\AttributeTok{A\_dass =} \DecValTok{74}\SpecialCharTok{:}\DecValTok{80}\NormalTok{ , }\AttributeTok{D\_dass =} \DecValTok{81}\SpecialCharTok{:}\DecValTok{87}\NormalTok{ , }\AttributeTok{A\_hads =} \DecValTok{88}\SpecialCharTok{:}\DecValTok{94}\NormalTok{, }\AttributeTok{D\_hads =} \DecValTok{95}\SpecialCharTok{:}\DecValTok{101}\NormalTok{, }\AttributeTok{A\_gad =} \DecValTok{102}\SpecialCharTok{:}\DecValTok{108}\NormalTok{, }\AttributeTok{A\_sas =} \DecValTok{109}\SpecialCharTok{:}\DecValTok{123}\NormalTok{)}

\CommentTok{\# To be more readable on Network Graphs here is a vector of shortenned names}

\NormalTok{readable\_anx\_dp }\OtherTok{\textless{}{-}} \FunctionTok{c}\NormalTok{(}\StringTok{"c01"}\NormalTok{, }\StringTok{"c02"}\NormalTok{, }\StringTok{"c03"}\NormalTok{, }\StringTok{"c04"}\NormalTok{, }\StringTok{"c05"}\NormalTok{, }\StringTok{"c06"}\NormalTok{, }\StringTok{"c07"}\NormalTok{, }\StringTok{"c08"}\NormalTok{, }\StringTok{"c09"}\NormalTok{, }\StringTok{"c10"}\NormalTok{, }\StringTok{"c11"}\NormalTok{, }\StringTok{"c12"}\NormalTok{, }\StringTok{"c13"}\NormalTok{, }\StringTok{"c14"}\NormalTok{, }\StringTok{"c15"}\NormalTok{, }\StringTok{"c16"}\NormalTok{, }\StringTok{"c17"}\NormalTok{, }\StringTok{"c18"}\NormalTok{, }\StringTok{"c19"}\NormalTok{, }\StringTok{"c20"}\NormalTok{, }
\StringTok{"ph01"}\NormalTok{, }\StringTok{"ph02"}\NormalTok{, }\StringTok{"ph03"}\NormalTok{, }\StringTok{"ph04"}\NormalTok{, }\StringTok{"ph05"}\NormalTok{, }\StringTok{"ph06"}\NormalTok{, }\StringTok{"ph07"}\NormalTok{, }\StringTok{"ph08"}\NormalTok{, }\StringTok{"ph09"}\NormalTok{, }
\StringTok{"pr01"}\NormalTok{, }\StringTok{"pr02"}\NormalTok{, }\StringTok{"pr03"}\NormalTok{, }\StringTok{"pr04"}\NormalTok{, }\StringTok{"pr05"}\NormalTok{, }\StringTok{"pr06"}\NormalTok{, }\StringTok{"pr07"}\NormalTok{, }\StringTok{"pr08"}\NormalTok{, }\StringTok{"pr09"}\NormalTok{, }\StringTok{"pr10"}\NormalTok{, }\StringTok{"pr11"}\NormalTok{, }\StringTok{"pr12"}\NormalTok{, }\StringTok{"pr13"}\NormalTok{, }\StringTok{"pr14"}\NormalTok{, }\StringTok{"pr15"}\NormalTok{, }\StringTok{"pr16"}\NormalTok{, }\StringTok{"pr17"}\NormalTok{, }\StringTok{"pr18"}\NormalTok{, }\StringTok{"pr19"}\NormalTok{, }\StringTok{"pr20"}\NormalTok{, }\StringTok{"pr21"}\NormalTok{, }\StringTok{"pr22"}\NormalTok{, }\StringTok{"pr23"}\NormalTok{, }\StringTok{"pr24"}\NormalTok{, }\StringTok{"pr25"}\NormalTok{, }\StringTok{"pr26"}\NormalTok{, }\StringTok{"pr27"}\NormalTok{, }\StringTok{"pr28"}\NormalTok{, }
\StringTok{"q01"}\NormalTok{, }\StringTok{"q02"}\NormalTok{, }\StringTok{"q03"}\NormalTok{, }\StringTok{"q04"}\NormalTok{, }\StringTok{"q05"}\NormalTok{, }\StringTok{"q06"}\NormalTok{, }\StringTok{"q07"}\NormalTok{, }\StringTok{"q08"}\NormalTok{, }\StringTok{"q09"}\NormalTok{, }\StringTok{"q10"}\NormalTok{, }\StringTok{"q11"}\NormalTok{, }\StringTok{"q12"}\NormalTok{, }\StringTok{"q13"}\NormalTok{, }\StringTok{"q14"}\NormalTok{, }\StringTok{"q15"}\NormalTok{, }\StringTok{"q16"}\NormalTok{, }
\StringTok{"d02"}\NormalTok{, }\StringTok{"d04"}\NormalTok{, }\StringTok{"d07"}\NormalTok{, }\StringTok{"d09"}\NormalTok{, }\StringTok{"d15"}\NormalTok{, }\StringTok{"d19"}\NormalTok{, }\StringTok{"d20"}\NormalTok{, }\StringTok{"d03"}\NormalTok{, }\StringTok{"d05"}\NormalTok{, }\StringTok{"d10"}\NormalTok{, }\StringTok{"d13"}\NormalTok{, }\StringTok{"d16"}\NormalTok{, }\StringTok{"d17"}\NormalTok{, }\StringTok{"d21"}\NormalTok{, }
\StringTok{"h01"}\NormalTok{, }\StringTok{"h03"}\NormalTok{, }\StringTok{"h05"}\NormalTok{, }\StringTok{"h07"}\NormalTok{, }\StringTok{"h09"}\NormalTok{, }\StringTok{"h11"}\NormalTok{, }\StringTok{"h13"}\NormalTok{, }\StringTok{"h02"}\NormalTok{, }\StringTok{"h04"}\NormalTok{, }\StringTok{"h06"}\NormalTok{, }\StringTok{"h08"}\NormalTok{, }\StringTok{"h10"}\NormalTok{, }\StringTok{"h12"}\NormalTok{, }\StringTok{"h14"}\NormalTok{, }
\StringTok{"g01"}\NormalTok{, }\StringTok{"g02"}\NormalTok{, }\StringTok{"g03"}\NormalTok{, }\StringTok{"g04"}\NormalTok{, }\StringTok{"g05"}\NormalTok{, }\StringTok{"g06"}\NormalTok{, }\StringTok{"g07"}\NormalTok{, }
\StringTok{"s01"}\NormalTok{, }\StringTok{"s02"}\NormalTok{, }\StringTok{"s03"}\NormalTok{, }\StringTok{"s04"}\NormalTok{, }\StringTok{"s05"}\NormalTok{, }\StringTok{"s06"}\NormalTok{, }\StringTok{"s07"}\NormalTok{, }\StringTok{"s08"}\NormalTok{, }\StringTok{"s09"}\NormalTok{, }\StringTok{"s10"}\NormalTok{, }\StringTok{"s11"}\NormalTok{, }\StringTok{"s12"}\NormalTok{, }\StringTok{"s13"}\NormalTok{, }\StringTok{"s14"}\NormalTok{, }\StringTok{"s15"}\NormalTok{)  }
\CommentTok{\#"t01", "t02", "t03", "t04", "t05", "t06", "t07", "t08", "t09", "t10", "t11", "t12", "t13", "t14", "t15", "t16", "t17", \#"t18", "t19", "t20", "t21", "t22", "t23", "t24", "t25", "t26", "t27", "t28", "t29", "t30", "t31", "t32", "t33", "t34", \#"t35", "t36", "t37", "t38", "t39", "t40", "t41", "t42", "t43", "t44", "t45", "t46", "t47", "t48", "t49", "t50")}

\CommentTok{\# Group names}
\NormalTok{group\_names\_anx\_dep }\OtherTok{\textless{}{-}} \FunctionTok{c}\NormalTok{(}\StringTok{"Depr\_cesd"}\NormalTok{, }\StringTok{"Depr\_phq"}\NormalTok{, }\StringTok{"Depr\_promis"}\NormalTok{, }\StringTok{"Depr\_qids"}\NormalTok{, }\StringTok{"Ans\_dass"}\NormalTok{ , }\StringTok{"Depr\_dass"}\NormalTok{ , }\StringTok{"Ans\_hads"}\NormalTok{, }\StringTok{"Depr\_hads"}\NormalTok{, }\StringTok{"Ans\_gad"}\NormalTok{, }\StringTok{"Ans\_sas"}\NormalTok{)}
\end{Highlighting}
\end{Shaded}

\subsection{Descriptive Statistics}\label{descriptive-statistics}

\subsubsection{Demographics}\label{demographics}

\begin{Shaded}
\begin{Highlighting}[]
\CommentTok{\#to be done}
\end{Highlighting}
\end{Shaded}

\subsubsection{General Statistics about the
scales}\label{general-statistics-about-the-scales}

Summary, Graphs, range, mean, sd, etc\ldots{}

\begin{Shaded}
\begin{Highlighting}[]
\CommentTok{\#to be done}
\end{Highlighting}
\end{Shaded}

\subsection{Force-directed plotting with
Fruchterman--Reingold}\label{force-directed-plotting-with-fruchtermanreingold}

Most studies on networks in psychopathology have used the
Fruchterman-Reingold (FR) algorithm to create graphs (Borsboom et al.,
2021). This algorithm is like setting up a system of balls connected by
stretchy strings. The strings pull connected balls closer, while the
other balls push them apart, resulting in a graph where nodes are
well-spaced and edges are evenly distributed (Deserno et al., 2022).

The goal of force-directed algorithms is not to place nodes in specific
meaningful locations, but to make the network easy to view, showing
connections and clusters clearly. The arrangement produced by these
algorithms is primarily driven by the need for visual clarity and does
not necessarily reflect any inherent structure or meaning beyond the
connections themselves. Therefore, it is crucial to focus on the
connections between nodes and their strengths, rather than the spatial
arrangement of nodes, when interpreting a force-directed graph (Deserno
et al., 2022).

The spatial positions of nodes in these graphs are highly sensitive to
variations in edge weights, even those as small as 0.01. A slight
alteration in edge weights can lead to significant changes in the node
arrangement (Burger et al., 2022). Therefore, the layout should not be
interpreted as reflecting inherent properties or relationships among the
nodes beyond their connection.

\subsubsection{Testing different correlation
methods}\label{testing-different-correlation-methods}

Spearman Correlation, as mentioned, is particularly useful for analyzing
non-parametric data. Spearman correlations, or rank-based correlations,
can be used to handle non-normal continuous or ordered categorical data.
On the other hand, automatic recognition of data type with cor\_auto
function from qgraph typically creates a dense correlation matrix,
incorporating automatic correction for ordinal and continuous data,
which can lead to multicollinearity. ForcePD forces the matrix to be
positive definite. A positive definite matrix guarantees that all its
eigenvalues are positive, a condition crucial for methods like testing
the robustness of the network. Multicollinearity arises when predictor
variables in a model are highly correlated, making it challenging to
discern the unique effect of each predictor (Epskamp, Borsboom, \&
Fried, 2018). Therefore, testing which adjencency matrix is better in
terms of multicollinearity and knowing whether it is positive. The
results below show that both matrices are positive definite but what
actually indicates that Spearman is the most suitable method for our
data is the amount of variable pairs over a threshold of 0.8. While
Spearman hasarount 34 pairs with that extremelly high correlation, the
cor\_auto function has 1232 pairs.

\begin{Shaded}
\begin{Highlighting}[]
\CommentTok{\# cor\_auto}
\NormalTok{symp\_auto }\OtherTok{\textless{}{-}} \FunctionTok{cor\_auto}\NormalTok{(df\_anx\_dp\_adj, }\AttributeTok{forcePD =} \ConstantTok{TRUE}\NormalTok{)}

\CommentTok{\# Spearman Correlation}

\NormalTok{symp\_spear }\OtherTok{\textless{}{-}} \FunctionTok{cor}\NormalTok{(df\_anx\_dp\_adj, }\AttributeTok{method =} \StringTok{"spearman"}\NormalTok{)}

\CommentTok{\# Estimating and penalizing connections}
\NormalTok{n1 }\OtherTok{\textless{}{-}} \FunctionTok{estimateNetwork}\NormalTok{(df\_anx\_dp\_adj, }\AttributeTok{default=}\StringTok{"EBICglasso"}\NormalTok{,  }\AttributeTok{corMethod =} \StringTok{"spearman"}\NormalTok{)}


\NormalTok{mc\_auto }\OtherTok{\textless{}{-}} \FunctionTok{multiC}\NormalTok{(symp\_auto)}
\FunctionTok{head}\NormalTok{(mc\_auto, }\DecValTok{5}\NormalTok{)}
\end{Highlighting}
\end{Shaded}

\begin{verbatim}
##   Variable1 Variable2 Correlation
## 1   cesd_06   cesd_04   0.9395364
## 2   cesd_04   cesd_06   0.9395364
## 3   cesd_04   cesd_02   0.9318807
## 4   cesd_02   cesd_04   0.9318807
## 5 promis_12 promis_10   0.9257129
\end{verbatim}

\begin{Shaded}
\begin{Highlighting}[]
\FunctionTok{cat}\NormalTok{(}\StringTok{"}\SpecialCharTok{\textbackslash{}n}\StringTok{Automatic correlation method is positive definite = "}\NormalTok{, }\FunctionTok{is.positive.definite}\NormalTok{(symp\_auto), }\StringTok{"}\SpecialCharTok{\textbackslash{}n}\StringTok{Its amount of variable pairs over the threshold of 0.8 is:"}\NormalTok{, }\FunctionTok{nrow}\NormalTok{(mc\_auto))}
\end{Highlighting}
\end{Shaded}

\begin{verbatim}
## 
## Automatic correlation method is positive definite =  TRUE 
## Its amount of variable pairs over the threshold of 0.8 is: 1300
\end{verbatim}

\begin{Shaded}
\begin{Highlighting}[]
\NormalTok{mc\_spear }\OtherTok{\textless{}{-}} \FunctionTok{multiC}\NormalTok{(symp\_spear)}
\FunctionTok{head}\NormalTok{(mc\_spear, }\DecValTok{5}\NormalTok{)}
\end{Highlighting}
\end{Shaded}

\begin{verbatim}
##   Variable1 Variable2 Correlation
## 1 promis_12 promis_10   0.8486060
## 2 promis_10 promis_12   0.8486060
## 3 promis_05 promis_02   0.8340622
## 4 promis_02 promis_05   0.8340622
## 5   cesd_06   cesd_04   0.8288046
\end{verbatim}

\begin{Shaded}
\begin{Highlighting}[]
\FunctionTok{cat}\NormalTok{(}\StringTok{"}\SpecialCharTok{\textbackslash{}n}\StringTok{Spearman method is positive definite = "}\NormalTok{, }\FunctionTok{is.positive.definite}\NormalTok{(symp\_spear), }\StringTok{"}\SpecialCharTok{\textbackslash{}n}\StringTok{Its amount of variable pairs over the threshold of 0.8 is:"}\NormalTok{, }\FunctionTok{nrow}\NormalTok{(mc\_spear))}
\end{Highlighting}
\end{Shaded}

\begin{verbatim}
## 
## Spearman method is positive definite =  TRUE 
## Its amount of variable pairs over the threshold of 0.8 is: 34
\end{verbatim}

\begin{Shaded}
\begin{Highlighting}[]
\NormalTok{mc\_n1 }\OtherTok{\textless{}{-}} \FunctionTok{multiC}\NormalTok{(n1}\SpecialCharTok{$}\NormalTok{graph)}
\FunctionTok{head}\NormalTok{(mc\_n1, }\DecValTok{5}\NormalTok{)}
\end{Highlighting}
\end{Shaded}

\begin{verbatim}
## [1] Variable1   Variable2   Correlation
## <0 rows> (or 0-length row.names)
\end{verbatim}

\begin{Shaded}
\begin{Highlighting}[]
\FunctionTok{cat}\NormalTok{(}\StringTok{"}\SpecialCharTok{\textbackslash{}n}\StringTok{Spearman method is positive definite = "}\NormalTok{, }\FunctionTok{is.positive.definite}\NormalTok{(n1}\SpecialCharTok{$}\NormalTok{graph), }\StringTok{"}\SpecialCharTok{\textbackslash{}n}\StringTok{Its amount of variable pairs over the threshold of 0.8 is:"}\NormalTok{, }\FunctionTok{nrow}\NormalTok{(mc\_n1))}
\end{Highlighting}
\end{Shaded}

\begin{verbatim}
## 
## Spearman method is positive definite =  FALSE 
## Its amount of variable pairs over the threshold of 0.8 is: 0
\end{verbatim}

\subsubsection{Creating networks out of individual scales using FR
algorithm and Spearman
Correlation}\label{creating-networks-out-of-individual-scales-using-fr-algorithm-and-spearman-correlation}

The Spearman correlation is particularly useful for analyzing
non-parametric data, such as those generated by Likert scales, which are
commonly employed in psychological research. Likert scales typically
consist of a set of statements or questions with ordered response
categories, such as ``strongly agree,'' and ``strongly disagree''. The
ordered nature of these responses often violates the assumptions of
parametric statistical methods that require continuous, normally
distributed data (Blanken et al., 2022)*. Spearman correlation, as a
non-parametric method, does not rely on these assumptions and can be
used to assess the strength and direction of the association between
variables measured on Likert scales. This makes it a valuable tool for
examining relationships between variables that do not meet the criteria
for parametric analyses, providing researchers with a robust alternative
for analyzing data commonly encountered in psychological research.

\begin{Shaded}
\begin{Highlighting}[]
\NormalTok{n }\OtherTok{\textless{}{-}} \DecValTok{1}
\ControlFlowTok{for}\NormalTok{(sc }\ControlFlowTok{in}\NormalTok{ groups\_anx\_dep\_adj)\{}
\NormalTok{  zero\_order\_scale }\OtherTok{\textless{}{-}} \FunctionTok{cor}\NormalTok{(df\_anx\_dp\_adj[sc], }\AttributeTok{method =} \StringTok{"spearman"}\NormalTok{)}
  \FunctionTok{qgraph}\NormalTok{(zero\_order\_scale, }\AttributeTok{layout=}\StringTok{"spring"}\NormalTok{, }\AttributeTok{labels =} \FunctionTok{colnames}\NormalTok{(zero\_order\_scale), }\AttributeTok{color =} \FunctionTok{sample}\NormalTok{(}\FunctionTok{colors}\NormalTok{(), }\DecValTok{1}\NormalTok{), }\AttributeTok{theme=}\StringTok{"TeamFortress"}\NormalTok{, }\AttributeTok{height =} \DecValTok{180}\NormalTok{, }\AttributeTok{width =} \DecValTok{140}\NormalTok{, }\AttributeTok{title =} \FunctionTok{paste}\NormalTok{(}\StringTok{"Network of"}\NormalTok{, group\_names\_anx\_dep[n]))}

\NormalTok{  n }\OtherTok{\textless{}{-}}\NormalTok{ n }\SpecialCharTok{+} \DecValTok{1}  
\NormalTok{\}}
\end{Highlighting}
\end{Shaded}

\includegraphics{network_analysis_files/figure-latex/unnamed-chunk-9-1.pdf}
\includegraphics{network_analysis_files/figure-latex/unnamed-chunk-9-2.pdf}
\includegraphics{network_analysis_files/figure-latex/unnamed-chunk-9-3.pdf}
\includegraphics{network_analysis_files/figure-latex/unnamed-chunk-9-4.pdf}
\includegraphics{network_analysis_files/figure-latex/unnamed-chunk-9-5.pdf}
\includegraphics{network_analysis_files/figure-latex/unnamed-chunk-9-6.pdf}
\includegraphics{network_analysis_files/figure-latex/unnamed-chunk-9-7.pdf}
\includegraphics{network_analysis_files/figure-latex/unnamed-chunk-9-8.pdf}
\includegraphics{network_analysis_files/figure-latex/unnamed-chunk-9-9.pdf}
\includegraphics{network_analysis_files/figure-latex/unnamed-chunk-9-10.pdf}

\subsection{Network graphs and
analysis}\label{network-graphs-and-analysis}

At first glace, the items seem to be rather clustered together. In spite
of the visual impression, it is crucial to explore some measures in
order to explore the network.

\begin{verbatim}
## pdf 
##   2
\end{verbatim}

\includegraphics{network_analysis_files/figure-latex/unnamed-chunk-10-1.pdf}

\textless\textless\textless\textless\textless\textless\textless{} HEAD
\#\#\# Networks by nosological entities aimed by the scale ======= \#\#
Centrality Measures for Individual Scales (Strength, Betweenness,
Closeness):

\begin{Shaded}
\begin{Highlighting}[]
\NormalTok{spearman\_cesd }\OtherTok{\textless{}{-}} \FunctionTok{cor}\NormalTok{(D\_cesd, }\AttributeTok{method =} \StringTok{"spearman"}\NormalTok{)}
\NormalTok{centrality\_cesd }\OtherTok{\textless{}{-}} \FunctionTok{centrality\_auto}\NormalTok{(spearman\_cesd)}

\NormalTok{spearman\_phq }\OtherTok{\textless{}{-}} \FunctionTok{cor}\NormalTok{(D\_phq, }\AttributeTok{method =} \StringTok{"spearman"}\NormalTok{)}
\NormalTok{centrality\_phq }\OtherTok{\textless{}{-}} \FunctionTok{centrality\_auto}\NormalTok{(spearman\_phq)}

\NormalTok{spearman\_promis }\OtherTok{\textless{}{-}} \FunctionTok{cor}\NormalTok{(D\_promis, }\AttributeTok{method =} \StringTok{"spearman"}\NormalTok{)}
\NormalTok{centrality\_promis }\OtherTok{\textless{}{-}} \FunctionTok{centrality\_auto}\NormalTok{(spearman\_promis)}

\NormalTok{spearman\_qids }\OtherTok{\textless{}{-}} \FunctionTok{cor}\NormalTok{(D\_qids, }\AttributeTok{method =} \StringTok{"spearman"}\NormalTok{)}
\NormalTok{centrality\_qids }\OtherTok{\textless{}{-}} \FunctionTok{centrality\_auto}\NormalTok{(spearman\_qids)}

\NormalTok{spearman\_dass\_a }\OtherTok{\textless{}{-}} \FunctionTok{cor}\NormalTok{(A\_dass, }\AttributeTok{method =} \StringTok{"spearman"}\NormalTok{)}
\NormalTok{centrality\_dass\_a }\OtherTok{\textless{}{-}} \FunctionTok{centrality\_auto}\NormalTok{(spearman\_dass\_a)}

\NormalTok{spearman\_dass\_d }\OtherTok{\textless{}{-}} \FunctionTok{cor}\NormalTok{(D\_dass, }\AttributeTok{method =} \StringTok{"spearman"}\NormalTok{)}
\NormalTok{centrality\_dass\_d }\OtherTok{\textless{}{-}} \FunctionTok{centrality\_auto}\NormalTok{(spearman\_dass\_d)}

\NormalTok{spearman\_hads\_a }\OtherTok{\textless{}{-}} \FunctionTok{cor}\NormalTok{(A\_hads, }\AttributeTok{method =} \StringTok{"spearman"}\NormalTok{)}
\NormalTok{centrality\_hads\_a }\OtherTok{\textless{}{-}} \FunctionTok{centrality\_auto}\NormalTok{(spearman\_hads\_a)}

\NormalTok{spearman\_hads\_d }\OtherTok{\textless{}{-}} \FunctionTok{cor}\NormalTok{(D\_hads, }\AttributeTok{method =} \StringTok{"spearman"}\NormalTok{)}
\NormalTok{centrality\_hads\_d }\OtherTok{\textless{}{-}} \FunctionTok{centrality\_auto}\NormalTok{(spearman\_hads\_d)}

\NormalTok{spearman\_gad }\OtherTok{\textless{}{-}} \FunctionTok{cor}\NormalTok{(A\_gad, }\AttributeTok{method =} \StringTok{"spearman"}\NormalTok{)}
\NormalTok{centrality\_gad }\OtherTok{\textless{}{-}} \FunctionTok{centrality\_auto}\NormalTok{(spearman\_gad)}

\NormalTok{spearman\_sas }\OtherTok{\textless{}{-}} \FunctionTok{cor}\NormalTok{(A\_sas, }\AttributeTok{method =} \StringTok{"spearman"}\NormalTok{)}
\NormalTok{centrality\_sas }\OtherTok{\textless{}{-}} \FunctionTok{centrality\_auto}\NormalTok{(spearman\_sas)}

\NormalTok{spearman\_tmas }\OtherTok{\textless{}{-}} \FunctionTok{cor}\NormalTok{(A\_tmas, }\AttributeTok{method =} \StringTok{"spearman"}\NormalTok{)}
\NormalTok{centrality\_tmas }\OtherTok{\textless{}{-}} \FunctionTok{centrality\_auto}\NormalTok{(spearman\_tmas)}
\end{Highlighting}
\end{Shaded}

\begin{quote}
\begin{quote}
\begin{quote}
\begin{quote}
\begin{quote}
\begin{quote}
\begin{quote}
d5b403185badaf959096b494ef3bf2a688a3cf61
\end{quote}
\end{quote}
\end{quote}
\end{quote}
\end{quote}
\end{quote}
\end{quote}

Identifying nodes only by the nosological entity one can see that the
quantity of items related to depressive symptoms is far superior. It is
no surprise since we excluded TMAS from the dataset.Anxiety items are in
total 36 and the other 87 are related to depression. Despite
representing a different nosological entity, these items seem to be
clustered together with those of depression. Later, we investigate if
the items about anxiety will form their own communities or blend with
the others.

\includegraphics{network_analysis_files/figure-latex/unnamed-chunk-12-1.pdf}

\subsection{4. Robustness Network}\label{robustness-network}

Once we have find a correlation matrix that is positive determinate is
controlled for multicollinearity, it is necessary to test the robustness
of this weighted network. In order to do it, we will use the bootnet
package (Version 1.6; Epskamp, 2024). Bootnet uses bootstrapping to
calculate standard errors and confidence intervals and is able to
poerform it in nonparametric data, such as ours.

\subsubsection{Edges Robustness}\label{edges-robustness}

In the graph below the red dots represent the mean from each edge of the
original data and the black dots represent the same but for the
bootstrap edges. In spite of the wider confidence interval band when it
comes close to zero, all the dots seem to be tight together, indicating
strong stability and reliability.

\begin{Shaded}
\begin{Highlighting}[]
\DocumentationTok{\#\#\# 4.1 edge weight bootstrap}

\CommentTok{\#m1\textless{}{-}bootnet(df\_anx\_dp\_adj, nBoots=500, default="EBICglasso", labels=colnames(df\_anx\_dp\_adj), type="nonparametric", statistics =c("edge", "strength", "closeness", "betweenness"))}
\end{Highlighting}
\end{Shaded}

\begin{Shaded}
\begin{Highlighting}[]
\CommentTok{\#save(m1, file = "boot\_symp\_df\_edgeweights.Rdata")}
\FunctionTok{load}\NormalTok{(}\AttributeTok{file =} \StringTok{"../images/boot\_symp\_df\_edgeweights.Rdata"}\NormalTok{)}

\FunctionTok{plot}\NormalTok{(m1, }\AttributeTok{labels =} \ConstantTok{FALSE}\NormalTok{, }\AttributeTok{order =} \StringTok{"sample"}\NormalTok{)}
\end{Highlighting}
\end{Shaded}

\includegraphics{network_analysis_files/figure-latex/unnamed-chunk-14-1.pdf}

\begin{Shaded}
\begin{Highlighting}[]
\CommentTok{\#pdf("../images/boot\_symp\_df\_edgeweights.pdf", useDingbats=FALSE) }
\CommentTok{\#plot(m1, labels = F, order = "sample")}
\CommentTok{\#dev.off()}
\end{Highlighting}
\end{Shaded}

\subsubsection{Node dropping}\label{node-dropping}

This bootstring method removes the nodes progressively and measures the
correlation of the network for each removal. When there is only around
35\% of the original nodes the correlation goes down quickly, increasing
variability. This suggests that the network is significantly stable
since it takes a huge amount of nodes removal to make it reduce the
reliability of general node strength.

\begin{Shaded}
\begin{Highlighting}[]
\CommentTok{\#m2\textless{}{-} bootnet(df\_anx\_dp\_adj, default = "EBICglasso", nBoots=500, type = "node", labels=colnames(df\_anx\_dp\_adj), statistics =c("strength", "closeness", "betweenness"))}
\end{Highlighting}
\end{Shaded}

\begin{Shaded}
\begin{Highlighting}[]
\CommentTok{\#save(m2, file = "boot\_symp\_df\_nodedropping.Rdata")}
\FunctionTok{load}\NormalTok{(}\AttributeTok{file =} \StringTok{"../images/boot\_symp\_df\_nodedropping.Rdata"}\NormalTok{)}

\FunctionTok{plot}\NormalTok{(m2)}
\end{Highlighting}
\end{Shaded}

\includegraphics{network_analysis_files/figure-latex/unnamed-chunk-16-1.pdf}

\begin{Shaded}
\begin{Highlighting}[]
\CommentTok{\#plot(m2, perNode = T, legend = T, statistics =c("strength"))}
\end{Highlighting}
\end{Shaded}

\subsection{Centrality Measures (Strength, Betweenness,
Closeness):}\label{centrality-measures-strength-betweenness-closeness}

Centrality measures are used to understand the relative importance or
influence of nodes within a network. It is important to acknowledge that
while centrality measures can provide valuable insights into the
structure of a network, they should not be directly interpreted as
measures of causal influence (Borsboom, 2017). The sources emphasize
that network models, especially those based on correlation or partial
correlation, do not inherently specify the direction of causal effects
(Borsboom et al., 2021).

These measures help identify the most important nodes (symptoms) in the
network. Strength indicates how connected a node is (Cai et al., 2022);
Betweenness represents how often a node lies on the shortest path
between other nodes; Closeness measures how quickly a node can reach
other nodes.

\begin{Shaded}
\begin{Highlighting}[]
\CommentTok{\# creating the table with all measures requested}
\NormalTok{cent\_table\_spear }\OtherTok{\textless{}{-}} \FunctionTok{centrality\_auto}\NormalTok{(symp\_spear)}

\DocumentationTok{\#\#\# Centrality Measures}
\FunctionTok{pdf}\NormalTok{(}\StringTok{"../images/centrality\_spear.pdf"}\NormalTok{)}

\NormalTok{cent\_graph\_spear }\OtherTok{\textless{}{-}} \FunctionTok{centralityPlot}\NormalTok{(symp\_spear, }\AttributeTok{scale =} \StringTok{"raw0"}\NormalTok{, }\AttributeTok{include =} \FunctionTok{c}\NormalTok{(}\StringTok{"Strength"}\NormalTok{, }\StringTok{"Closeness"}\NormalTok{, }\StringTok{"Betweenness"}\NormalTok{), }\AttributeTok{decreasing =}\NormalTok{ T, }\AttributeTok{labels =}\NormalTok{ readable\_anx\_dp)}
\FunctionTok{dev.off}\NormalTok{()}
\end{Highlighting}
\end{Shaded}

\begin{verbatim}
## pdf 
##   2
\end{verbatim}

\begin{Shaded}
\begin{Highlighting}[]
\NormalTok{cent\_graph\_spear}
\end{Highlighting}
\end{Shaded}

\includegraphics{network_analysis_files/figure-latex/unnamed-chunk-17-1.pdf}

\#\#\#\#Strength

Strength is the sum of the correlations of one node to all other nodes
with higher values reflecting greater centrality in the network (Cai et
al., 2022).

There is a remarkable dominance of promis items among the strongest
ones. Almost all of the top 10 are arguably items related to the core
DSM-V-TR symptom ``depressed mood''. The data distribution for
``Strength'' is substantially skewed to the right, as presented in the
graph. Most items have moderate to high strength, which could indicate
that many nodes are strong within their own clusters but less connected
across clusters. This hypothesis makes even more sense when we take into
consideration the near-zero correlation between Strength and Betweenness
measures (-0.0187) that suggests that nodes with high strength are not
acting as bridges between clusters, but are tightely connected within
their groups. Core Major Depression symptoms like depressed mood and
anhedonia are particularly prominent here.

\begin{Shaded}
\begin{Highlighting}[]
\NormalTok{top\_str }\OtherTok{\textless{}{-}}\NormalTok{ cent\_table\_spear}\SpecialCharTok{$}\NormalTok{node.centrality }\SpecialCharTok{\%\textgreater{}\%}
\FunctionTok{arrange}\NormalTok{(}\FunctionTok{desc}\NormalTok{(Strength)) }

\FunctionTok{head}\NormalTok{(top\_str, }\DecValTok{10}\NormalTok{)}
\end{Highlighting}
\end{Shaded}

\begin{verbatim}
##           Betweenness   Closeness Strength ExpectedInfluence
## promis_05           1 0.004365505 72.76541          72.10187
## promis_21           4 0.004315606 71.70724          71.02427
## promis_28           2 0.004349012 71.62770          70.87923
## promis_15           0 0.004311767 71.46486          70.77000
## promis_02           0 0.004246198 70.95327          70.28330
## promis_24           0 0.004251037 70.55769          69.92529
## promis_10           0 0.004221163 70.51280          69.82857
## promis_07           0 0.004234416 70.41204          69.67582
## promis_03           2 0.004262541 70.38916          69.75880
## promis_26           1 0.004227029 70.03546          69.46545
\end{verbatim}

\begin{Shaded}
\begin{Highlighting}[]
\FunctionTok{tail}\NormalTok{(top\_str, }\DecValTok{10}\NormalTok{)}
\end{Highlighting}
\end{Shaded}

\begin{verbatim}
##         Betweenness   Closeness Strength ExpectedInfluence
## hads_11           0 0.002340342 36.32468         -36.32468
## qids_06           7 0.002086622 31.89005          30.88587
## hads_09           0 0.001943953 29.45306          28.99017
## qids_03           0 0.001790776 26.86590          26.61916
## qids_09           1 0.001643137 24.06010          22.96374
## cesd_11           0 0.001720255 23.87546          23.63657
## qids_07           0 0.001683269 21.98051          21.00595
## qids_08           0 0.001535180 21.80473          20.94286
## qids_02           0 0.001517898 18.92350          18.57312
## qids_04           0 0.001251495 14.29942          14.14807
\end{verbatim}

\begin{Shaded}
\begin{Highlighting}[]
\FunctionTok{cat}\NormalTok{(}\StringTok{"Top 10 Strength }\SpecialCharTok{\textbackslash{}n\textbackslash{}n}\StringTok{"}\NormalTok{)}
\end{Highlighting}
\end{Shaded}

\begin{verbatim}
## Top 10 Strength
\end{verbatim}

\begin{Shaded}
\begin{Highlighting}[]
\CommentTok{\#questions[c("promis\_05", "promis\_21", "promis\_28","promis\_15", "promis\_02")] }
\NormalTok{questions[}\FunctionTok{c}\NormalTok{(}\FunctionTok{which}\NormalTok{(}\FunctionTok{colnames}\NormalTok{(df\_anx\_dp\_adj) }\SpecialCharTok{\%in\%} \FunctionTok{sort}\NormalTok{(}\FunctionTok{row.names}\NormalTok{(top\_str)[}\DecValTok{1}\SpecialCharTok{:}\DecValTok{10}\NormalTok{])))] }
\end{Highlighting}
\end{Shaded}

\begin{verbatim}
## $promis_02
## [1] "Sentí que nada me ilusionaba"
## 
## $promis_03
## [1] "Me sentí indefenso/a (que no podía hacer nada para ayudarme)"
## 
## $promis_05
## [1] "Sentí que nada me podía animar"
## 
## $promis_07
## [1] "Me sentí triste"
## 
## $promis_10
## [1] "Me sentí fracasado/a"
## 
## $promis_15
## [1] "Me sentí deprimido/a"
## 
## $promis_21
## [1] "Me sentí desesperanzado/a"
## 
## $promis_24
## [1] "Sentí que nada me interesaba"
## 
## $promis_26
## [1] "Sentí que mi vida estaba vacía"
## 
## $promis_28
## [1] "Me sentí completamente agotado/a emocionalmente"
\end{verbatim}

\begin{Shaded}
\begin{Highlighting}[]
\FunctionTok{cat}\NormalTok{(}\StringTok{"}\SpecialCharTok{\textbackslash{}n}\StringTok{Bottom 20 Strength }\SpecialCharTok{\textbackslash{}n\textbackslash{}n}\StringTok{"}\NormalTok{)}
\end{Highlighting}
\end{Shaded}

\begin{verbatim}
## 
## Bottom 20 Strength
\end{verbatim}

\begin{Shaded}
\begin{Highlighting}[]
\CommentTok{\#questions[c("cesd\_11", "qids\_07", "qids\_08","qids\_02", "qids\_04")] }
\NormalTok{questions[}\FunctionTok{c}\NormalTok{(}\FunctionTok{which}\NormalTok{(}\FunctionTok{colnames}\NormalTok{(df\_anx\_dp\_adj) }\SpecialCharTok{\%in\%} \FunctionTok{sort}\NormalTok{(}\FunctionTok{row.names}\NormalTok{(top\_str)[}\DecValTok{104}\SpecialCharTok{:}\DecValTok{123}\NormalTok{])))] }
\end{Highlighting}
\end{Shaded}

\begin{verbatim}
## $cesd_01
## [1] "Tenía poco apetito"
## 
## $cesd_09
## [1] "Sentía que era una mala persona"
## 
## $cesd_11
## [1] "Dormía más de lo habitual"
## 
## $cesd_15
## [1] "Quería hacerse daño"
## 
## $cesd_18
## [1] "Perdió peso sin intentarlo"
## 
## $qids_01
## [1] "Problemas a la hora de dormirse"
## 
## $qids_02
## [1] "Sueño nocturno alterado"
## 
## $qids_03
## [1] "Despertarse demasiado temprano"
## 
## $qids_04
## [1] "Dormir demasiado"
## 
## $qids_06
## [1] "Ha habido disminución del apetito"
## 
## $qids_07
## [1] "Ha habido aumento del apetito"
## 
## $qids_08
## [1] "Ha habido pérdida de peso"
## 
## $qids_09
## [1] "Ha habido aumento de peso"
## 
## $hads_09
## [1] "Tengo una sensación extraña, como si tuviera mariposas en el estómago"
## 
## $hads_11
## [1] "Me siento inquieto/a, como si no pudiera parar de moverme"
## 
## $hads_14
## [1] "Me divierto con un buen libro, la radio, o un programa de televisión"
## 
## $sas_06
## [1] "Me molestan los dolores de cabeza, cuello y dolor de espalda"
## 
## $sas_10
## [1] "Tengo desmayos o sensación de desmayos"
## 
## $sas_13
## [1] "Tengo que vaciar mi vejiga con frecuencia"
## 
## $sas_14
## [1] "Se me calienta la cara y se me pone enrojecida"
\end{verbatim}

\begin{Shaded}
\begin{Highlighting}[]
\FunctionTok{summary}\NormalTok{(top\_str}\SpecialCharTok{$}\NormalTok{Strength)}
\end{Highlighting}
\end{Shaded}

\begin{verbatim}
##    Min. 1st Qu.  Median    Mean 3rd Qu.    Max. 
##   14.30   53.04   61.14   57.91   66.97   72.77
\end{verbatim}

\begin{Shaded}
\begin{Highlighting}[]
\NormalTok{top\_str}
\end{Highlighting}
\end{Shaded}

\begin{verbatim}
##           Betweenness   Closeness Strength ExpectedInfluence
## promis_05           1 0.004365505 72.76541          72.10187
## promis_21           4 0.004315606 71.70724          71.02427
## promis_28           2 0.004349012 71.62770          70.87923
## promis_15           0 0.004311767 71.46486          70.77000
## promis_02           0 0.004246198 70.95327          70.28330
## promis_24           0 0.004251037 70.55769          69.92529
## promis_10           0 0.004221163 70.51280          69.82857
## promis_07           0 0.004234416 70.41204          69.67582
## promis_03           2 0.004262541 70.38916          69.75880
## promis_26           1 0.004227029 70.03546          69.46545
## promis_16           0 0.004247373 69.84445          69.13925
## promis_23           0 0.004226995 69.80930          69.09570
## phq_02              0 0.004199282 69.74711          69.13069
## promis_19           0 0.004167235 69.62472          68.82916
## promis_25           0 0.004127187 69.51043          68.85726
## promis_12           0 0.004179470 69.50268          68.82247
## dass_13             0 0.004174714 69.34869          68.64348
## qids_05             3 0.004226433 69.22826          68.53338
## promis_13           0 0.004157139 69.18391          68.58634
## phq_06              0 0.004135586 68.71430          68.13412
## sas_04              7 0.004194146 68.41047          67.74584
## promis_04           0 0.004095054 67.73403          67.11397
## promis_17           8 0.004043714 67.71474          66.99984
## promis_01           3 0.004096780 67.55858          66.98393
## promis_06           0 0.004079000 67.49035          66.88234
## gad_04              7 0.004087235 67.36120          66.45176
## promis_11           0 0.004061058 67.31300          66.70663
## promis_18           0 0.004065504 67.13045          66.52045
## cesd_02             0 0.004095002 67.11337          66.46290
## dass_10             0 0.004043313 67.04027          66.47720
## cesd_17             0 0.004055556 66.98225          66.40246
## cesd_04             0 0.004066901 66.95792          66.34563
## promis_14           0 0.004008120 66.89616          66.28869
## promis_08           0 0.004050334 66.81481          66.26287
## cesd_06             0 0.004048560 66.73220          66.08963
## gad_02              0 0.003988842 65.85789          65.02218
## promis_27           2 0.004013408 65.81898          65.21915
## cesd_20             0 0.004021614 65.49416          64.81890
## promis_09           0 0.003942218 65.35468          64.71848
## cesd_08             4 0.003967253 65.35340          64.79722
## dass_16             6 0.003987713 65.20574          64.65450
## phq_01              0 0.003961423 65.03949          64.37323
## promis_22           0 0.003919588 64.84493          64.25700
## dass_21             3 0.003919996 64.53676          64.08782
## gad_07              1 0.003958294 64.26399          63.56580
## gad_03              0 0.003865575 63.93473          63.21244
## gad_05              2 0.003947502 63.83019          62.98978
## gad_01              4 0.003900918 63.68041          62.74014
## dass_17             0 0.003869430 63.66582          63.17903
## dass_15             0 0.003909181 63.45704          62.94351
## promis_20           0 0.003851671 62.95447          62.52618
## phq_07              0 0.003877604 62.92737          62.30008
## cesd_10             1 0.003864541 62.77519          62.16929
## gad_06              0 0.003826367 62.67095          61.95579
## sas_07              0 0.003850001 62.65721          61.94906
## dass_03             0 0.003856149 62.53519          61.90497
## sas_02              1 0.003886191 62.44665          61.71195
## dass_05             0 0.003807290 62.30481          61.68498
## sas_01              0 0.003821499 62.23201          61.35565
## sas_03              0 0.003814324 61.85838          61.14196
## cesd_16             0 0.003762756 61.17733          60.53834
## phq_04              0 0.003753760 61.14090          60.48965
## dass_20             0 0.003745179 60.99948          60.41252
## dass_09             0 0.003764482 60.99858          60.34265
## hads_13            16 0.003799790 60.91073          60.23454
## hads_03             0 0.003749486 60.70202          59.98904
## hads_01             0 0.003722996 60.56362          59.77613
## cesd_07             0 0.003707139 60.37425          59.80822
## qids_13             0 0.003684406 59.67399          59.12208
## cesd_03             0 0.003677162 59.57916          58.93214
## qids_11             7 0.003680638 59.53475          58.94672
## qids_14            13 0.003707683 59.04593          58.36471
## qids_10             0 0.003684553 58.74389          57.95845
## hads_05             0 0.003579597 58.65758          57.92425
## hads_02             0 0.003592106 58.49023          57.88499
## cesd_13             0 0.003634867 58.08730          57.42054
## hads_04             0 0.003562676 57.90110          57.31156
## qids_16             6 0.003582511 57.66363          56.70060
## dass_04             6 0.003628298 57.60559          57.07628
## phq_05             82 0.003644281 57.56202          56.90769
## dass_19             0 0.003592053 57.47805          56.87018
## hads_06             0 0.003496237 57.34211          56.73716
## phq_08              0 0.003603120 57.13730          56.64477
## qids_12             2 0.003515792 56.52034          56.16213
## hads_07             2 0.003489004 56.01847          55.29853
## hads_12             0 0.003440434 55.88313          55.28843
## hads_08             3 0.003499152 55.00555          54.34730
## phq_09            105 0.003422915 54.87306          54.54851
## sas_08              0 0.003433762 54.57348          53.91405
## cesd_14             0 0.003380698 54.45320          54.10374
## cesd_05            87 0.003403429 54.43822          53.76546
## qids_15             1 0.003388400 53.09728          52.53227
## sas_05              2 0.003331664 52.97355          52.50780
## phq_03             12 0.003286658 52.61150          51.98602
## dass_02             0 0.003301581 52.46602          51.81014
## sas_15              0 0.003296275 52.13075          51.58019
## cesd_12            53 0.003296921 51.77073          51.20715
## hads_10             1 0.003241431 51.72416          51.11984
## sas_09              0 0.003224010 50.73023          50.17564
## sas_11              0 0.003217828 50.70392          50.04507
## sas_12              0 0.003207713 50.50690          49.90328
## dass_07             0 0.003193651 50.43371          50.04382
## cesd_19             5 0.003148846 49.87952          49.23213
## hads_14             0 0.003107213 49.33823          48.77632
## cesd_09             0 0.003048508 48.49581          48.11945
## sas_06              0 0.002900075 46.21623          45.53144
## sas_10              0 0.002943170 46.13589          45.73820
## cesd_15             0 0.002867918 45.30223          45.03493
## sas_14              0 0.002851414 44.58144          44.08725
## cesd_01             2 0.002777187 42.74144          42.16380
## sas_13              0 0.002634876 41.49381          40.92050
## qids_01             0 0.002472004 38.55671          38.03457
## cesd_18            54 0.002465097 38.02216          37.62203
## hads_11             0 0.002340342 36.32468         -36.32468
## qids_06             7 0.002086622 31.89005          30.88587
## hads_09             0 0.001943953 29.45306          28.99017
## qids_03             0 0.001790776 26.86590          26.61916
## qids_09             1 0.001643137 24.06010          22.96374
## cesd_11             0 0.001720255 23.87546          23.63657
## qids_07             0 0.001683269 21.98051          21.00595
## qids_08             0 0.001535180 21.80473          20.94286
## qids_02             0 0.001517898 18.92350          18.57312
## qids_04             0 0.001251495 14.29942          14.14807
\end{verbatim}

\begin{Shaded}
\begin{Highlighting}[]
\FunctionTok{ggplot}\NormalTok{(top\_str, }\FunctionTok{aes}\NormalTok{(}\AttributeTok{x =}\NormalTok{ Strength)) }\SpecialCharTok{+}
  \FunctionTok{geom\_area}\NormalTok{(}\AttributeTok{stat =} \StringTok{"bin"}\NormalTok{, }\AttributeTok{bins =} \DecValTok{30}\NormalTok{, }\AttributeTok{fill =} \StringTok{"skyblue"}\NormalTok{, }\AttributeTok{alpha =} \FloatTok{0.6}\NormalTok{) }\SpecialCharTok{+}
  \FunctionTok{labs}\NormalTok{(}\AttributeTok{title =} \StringTok{"Distribution of Strength Values"}\NormalTok{,}
       \AttributeTok{x =} \StringTok{"Strength"}\NormalTok{,}
       \AttributeTok{y =} \StringTok{"Frequency"}\NormalTok{) }\SpecialCharTok{+}
  \FunctionTok{theme\_minimal}\NormalTok{()}
\end{Highlighting}
\end{Shaded}

\includegraphics{network_analysis_files/figure-latex/unnamed-chunk-18-1.pdf}

\#\#\#\#Betweenness

A node with high betweenness is a node that tends to connect clusters of
nodes together (Deserno et al., 2022). In our investigation, there is a
significant difference between the top 5 and all the rest, which
indicates that only certain items, and therefore only certain symptoms,
act as bridges between the different clusters. These items indicate the
constructs of weight/appetite issues, sleep issues and suicidal
thoughts. They essentially act as bridges between the most central nodes
and the peripherical nodes, which, in turn are have an importantantly
low value for closenness. In our case, the peripheral items (low
closenness) are connected to the strong central items through the ones
having high betweenness values.

\begin{Shaded}
\begin{Highlighting}[]
\NormalTok{top\_bet }\OtherTok{\textless{}{-}}\NormalTok{cent\_table\_spear}\SpecialCharTok{$}\NormalTok{node.centrality }\SpecialCharTok{\%\textgreater{}\%}
\FunctionTok{arrange}\NormalTok{(}\FunctionTok{desc}\NormalTok{(Betweenness)) }

\FunctionTok{head}\NormalTok{(top\_bet, }\DecValTok{10}\NormalTok{)}
\end{Highlighting}
\end{Shaded}

\begin{verbatim}
##           Betweenness   Closeness Strength ExpectedInfluence
## phq_09            105 0.003422915 54.87306          54.54851
## cesd_05            87 0.003403429 54.43822          53.76546
## phq_05             82 0.003644281 57.56202          56.90769
## cesd_18            54 0.002465097 38.02216          37.62203
## cesd_12            53 0.003296921 51.77073          51.20715
## hads_13            16 0.003799790 60.91073          60.23454
## qids_14            13 0.003707683 59.04593          58.36471
## phq_03             12 0.003286658 52.61150          51.98602
## promis_17           8 0.004043714 67.71474          66.99984
## qids_06             7 0.002086622 31.89005          30.88587
\end{verbatim}

\begin{Shaded}
\begin{Highlighting}[]
\CommentTok{\#tail(top\_bet, 10)}

\CommentTok{\#head(cent\_table\_spear$node.centrality \%\textgreater{}\%arrange(Betweenness, Strength), 10)}

\FunctionTok{cat}\NormalTok{(}\StringTok{"Top 10 Betweenness }\SpecialCharTok{\textbackslash{}n\textbackslash{}n}\StringTok{"}\NormalTok{)}
\end{Highlighting}
\end{Shaded}

\begin{verbatim}
## Top 10 Betweenness
\end{verbatim}

\begin{Shaded}
\begin{Highlighting}[]
\NormalTok{questions[}\FunctionTok{c}\NormalTok{(}\FunctionTok{which}\NormalTok{(}\FunctionTok{colnames}\NormalTok{(df\_anx\_dp\_adj) }\SpecialCharTok{\%in\%} \FunctionTok{sort}\NormalTok{(}\FunctionTok{row.names}\NormalTok{(top\_bet)[}\DecValTok{1}\SpecialCharTok{:}\DecValTok{10}\NormalTok{])))] }
\end{Highlighting}
\end{Shaded}

\begin{verbatim}
## $cesd_05
## [1] "Su sueño era inquieto"
## 
## $cesd_12
## [1] "Sentía que se movía muy lento"
## 
## $cesd_18
## [1] "Perdió peso sin intentarlo"
## 
## $phq_03
## [1] "Tener problemas para dormir o mantenerse dormido, o dormir demasiado"
## 
## $phq_05
## [1] "Tener poco apetito o comer en exceso"
## 
## $phq_09
## [1] "Tener pensamientos sobre estar muerto(a) o sobre hacerse daño a sí mismo(a) de alguna manera"
## 
## $promis_17
## [1] "Me sentí desanimado/a acerca del futuro"
## 
## $qids_06
## [1] "Ha habido disminución del apetito"
## 
## $qids_14
## [1] "Nivel de energía"
## 
## $hads_13
## [1] "Me asaltan sentimientos repentinos de pánico"
\end{verbatim}

\begin{Shaded}
\begin{Highlighting}[]
\FunctionTok{cat}\NormalTok{(}\StringTok{"}\SpecialCharTok{\textbackslash{}n}\StringTok{Bottom Betweenness: 83 items had betweenness = 0 }\SpecialCharTok{\textbackslash{}n\textbackslash{}n}\StringTok{"}\NormalTok{)}
\end{Highlighting}
\end{Shaded}

\begin{verbatim}
## 
## Bottom Betweenness: 83 items had betweenness = 0
\end{verbatim}

\begin{Shaded}
\begin{Highlighting}[]
\CommentTok{\#which(colnames(df\_anx\_dp\_adj) \%in\% sort(row.names(top\_bet)[1:40]))}
\end{Highlighting}
\end{Shaded}

\#\#\#\#Closeness

Closeness quantifies how far removed a node is from every other node in
the network (Deserno et al., 2022). As mentioned above, low closenness
means being more peripheral. In our network, given that the
intercorrelation between strength and closenness is almost perfect, if a
node is strong it is within a central cluster, well connected to the
other local nodes.

\begin{Shaded}
\begin{Highlighting}[]
\NormalTok{top\_clo }\OtherTok{\textless{}{-}}\NormalTok{ cent\_table\_spear}\SpecialCharTok{$}\NormalTok{node.centrality }\SpecialCharTok{\%\textgreater{}\%}
\FunctionTok{arrange}\NormalTok{(}\FunctionTok{desc}\NormalTok{(Closeness)) }

\NormalTok{bottom\_clo }\OtherTok{\textless{}{-}}\NormalTok{ cent\_table\_spear}\SpecialCharTok{$}\NormalTok{node.centrality }\SpecialCharTok{\%\textgreater{}\%}
\FunctionTok{arrange}\NormalTok{(Closeness)}

\FunctionTok{head}\NormalTok{(top\_clo, }\DecValTok{10}\NormalTok{) }
\end{Highlighting}
\end{Shaded}

\begin{verbatim}
##           Betweenness   Closeness Strength ExpectedInfluence
## promis_05           1 0.004365505 72.76541          72.10187
## promis_28           2 0.004349012 71.62770          70.87923
## promis_21           4 0.004315606 71.70724          71.02427
## promis_15           0 0.004311767 71.46486          70.77000
## promis_03           2 0.004262541 70.38916          69.75880
## promis_24           0 0.004251037 70.55769          69.92529
## promis_16           0 0.004247373 69.84445          69.13925
## promis_02           0 0.004246198 70.95327          70.28330
## promis_07           0 0.004234416 70.41204          69.67582
## promis_26           1 0.004227029 70.03546          69.46545
\end{verbatim}

\begin{Shaded}
\begin{Highlighting}[]
\FunctionTok{tail}\NormalTok{(top\_clo, }\DecValTok{10}\NormalTok{)}
\end{Highlighting}
\end{Shaded}

\begin{verbatim}
##         Betweenness   Closeness Strength ExpectedInfluence
## hads_11           0 0.002340342 36.32468         -36.32468
## qids_06           7 0.002086622 31.89005          30.88587
## hads_09           0 0.001943953 29.45306          28.99017
## qids_03           0 0.001790776 26.86590          26.61916
## cesd_11           0 0.001720255 23.87546          23.63657
## qids_07           0 0.001683269 21.98051          21.00595
## qids_09           1 0.001643137 24.06010          22.96374
## qids_08           0 0.001535180 21.80473          20.94286
## qids_02           0 0.001517898 18.92350          18.57312
## qids_04           0 0.001251495 14.29942          14.14807
\end{verbatim}

\begin{Shaded}
\begin{Highlighting}[]
\FunctionTok{cat}\NormalTok{(}\StringTok{"Top 10 Closenness }\SpecialCharTok{\textbackslash{}n\textbackslash{}n}\StringTok{"}\NormalTok{)}
\end{Highlighting}
\end{Shaded}

\begin{verbatim}
## Top 10 Closenness
\end{verbatim}

\begin{Shaded}
\begin{Highlighting}[]
\CommentTok{\#questions[c("promis\_05", "promis\_28", "promis\_21","promis\_15", "promis\_03")]}
\NormalTok{questions[}\FunctionTok{c}\NormalTok{(}\FunctionTok{which}\NormalTok{(}\FunctionTok{colnames}\NormalTok{(df\_anx\_dp\_adj) }\SpecialCharTok{\%in\%} \FunctionTok{sort}\NormalTok{(}\FunctionTok{row.names}\NormalTok{(top\_clo)[}\DecValTok{1}\SpecialCharTok{:}\DecValTok{10}\NormalTok{])))] }
\end{Highlighting}
\end{Shaded}

\begin{verbatim}
## $promis_02
## [1] "Sentí que nada me ilusionaba"
## 
## $promis_03
## [1] "Me sentí indefenso/a (que no podía hacer nada para ayudarme)"
## 
## $promis_05
## [1] "Sentí que nada me podía animar"
## 
## $promis_07
## [1] "Me sentí triste"
## 
## $promis_15
## [1] "Me sentí deprimido/a"
## 
## $promis_16
## [1] "Tuve problemas para tomar decisiones"
## 
## $promis_21
## [1] "Me sentí desesperanzado/a"
## 
## $promis_24
## [1] "Sentí que nada me interesaba"
## 
## $promis_26
## [1] "Sentí que mi vida estaba vacía"
## 
## $promis_28
## [1] "Me sentí completamente agotado/a emocionalmente"
\end{verbatim}

\begin{Shaded}
\begin{Highlighting}[]
\FunctionTok{cat}\NormalTok{(}\StringTok{"}\SpecialCharTok{\textbackslash{}n}\StringTok{Bottom 10 Closenness }\SpecialCharTok{\textbackslash{}n\textbackslash{}n}\StringTok{"}\NormalTok{)}
\end{Highlighting}
\end{Shaded}

\begin{verbatim}
## 
## Bottom 10 Closenness
\end{verbatim}

\begin{Shaded}
\begin{Highlighting}[]
\NormalTok{questions[}\FunctionTok{c}\NormalTok{(}\FunctionTok{which}\NormalTok{(}\FunctionTok{colnames}\NormalTok{(df\_anx\_dp\_adj) }\SpecialCharTok{\%in\%} \FunctionTok{sort}\NormalTok{(}\FunctionTok{row.names}\NormalTok{(bottom\_clo)[}\DecValTok{1}\SpecialCharTok{:}\DecValTok{10}\NormalTok{])))] }
\end{Highlighting}
\end{Shaded}

\begin{verbatim}
## $cesd_11
## [1] "Dormía más de lo habitual"
## 
## $qids_02
## [1] "Sueño nocturno alterado"
## 
## $qids_03
## [1] "Despertarse demasiado temprano"
## 
## $qids_04
## [1] "Dormir demasiado"
## 
## $qids_06
## [1] "Ha habido disminución del apetito"
## 
## $qids_07
## [1] "Ha habido aumento del apetito"
## 
## $qids_08
## [1] "Ha habido pérdida de peso"
## 
## $qids_09
## [1] "Ha habido aumento de peso"
## 
## $hads_09
## [1] "Tengo una sensación extraña, como si tuviera mariposas en el estómago"
## 
## $hads_11
## [1] "Me siento inquieto/a, como si no pudiera parar de moverme"
\end{verbatim}

\paragraph{Intercorrelation of Centrality
Measures:}\label{intercorrelation-of-centrality-measures}

High intercorrelations suggest that certain centrality measures may be
capturing similar network properties and these intercorrelations are
loud and clear: Strength and Closenness are intercorrelated almost
perfectly, while Betweenness has nearly no correlation with them.

\begin{Shaded}
\begin{Highlighting}[]
\FunctionTok{cat}\NormalTok{(}\StringTok{"Correlation of the measures}\SpecialCharTok{\textbackslash{}n\textbackslash{}n}\StringTok{"}\NormalTok{)}
\end{Highlighting}
\end{Shaded}

\begin{verbatim}
## Correlation of the measures
\end{verbatim}

\begin{Shaded}
\begin{Highlighting}[]
\FunctionTok{cat}\NormalTok{(}\StringTok{"Strength x Closeness: "}\NormalTok{, }\FunctionTok{cor}\NormalTok{(cent\_table\_spear}\SpecialCharTok{$}\NormalTok{node.centrality}\SpecialCharTok{$}\NormalTok{Strength, cent\_table\_spear}\SpecialCharTok{$}\NormalTok{node.centrality}\SpecialCharTok{$}\NormalTok{Closeness, }\AttributeTok{method=}\StringTok{"spearman"}\NormalTok{), }\StringTok{"}\SpecialCharTok{\textbackslash{}n}\StringTok{"}\NormalTok{)}
\end{Highlighting}
\end{Shaded}

\begin{verbatim}
## Strength x Closeness:  0.9969303
\end{verbatim}

\begin{Shaded}
\begin{Highlighting}[]
\FunctionTok{cat}\NormalTok{(}\StringTok{"Strength x Betweenness: "}\NormalTok{,}\FunctionTok{cor}\NormalTok{(cent\_table\_spear}\SpecialCharTok{$}\NormalTok{node.centrality}\SpecialCharTok{$}\NormalTok{Strength, cent\_table\_spear}\SpecialCharTok{$}\NormalTok{node.centrality}\SpecialCharTok{$}\NormalTok{Betweenness, }\AttributeTok{method=}\StringTok{"spearman"}\NormalTok{), }\StringTok{"}\SpecialCharTok{\textbackslash{}n}\StringTok{"}\NormalTok{)}
\end{Highlighting}
\end{Shaded}

\begin{verbatim}
## Strength x Betweenness:  -0.01877488
\end{verbatim}

\begin{Shaded}
\begin{Highlighting}[]
\FunctionTok{cat}\NormalTok{(}\StringTok{"Betweenness x Closeness: "}\NormalTok{,}\FunctionTok{cor}\NormalTok{(cent\_table\_spear}\SpecialCharTok{$}\NormalTok{node.centrality}\SpecialCharTok{$}\NormalTok{Betweenness, cent\_table\_spear}\SpecialCharTok{$}\NormalTok{node.centrality}\SpecialCharTok{$}\NormalTok{Closeness, }\AttributeTok{method=}\StringTok{"spearman"}\NormalTok{), }\StringTok{"}\SpecialCharTok{\textbackslash{}n}\StringTok{"}\NormalTok{)}
\end{Highlighting}
\end{Shaded}

\begin{verbatim}
## Betweenness x Closeness:  2.325542e-05
\end{verbatim}

Joining closenness and strength. There is no meaning in creating a
bottom 5 for betweenness since there are 83 nodes (67\%) having zero as
value for this measure.Instead, we will merge strength and closenness
only for visualization purposes and explore the 15 top and bottom
closenness/strength nodes. Also, we will visualize all the nodes that
have a betweenness value bigger than zero.

\includegraphics{network_analysis_files/figure-latex/unnamed-chunk-22-1.pdf}

\subsubsection{Clustering}\label{clustering}

We have been talking a lot about clustering, so it is time to start
investigating this feature. The clustering coefficient measures how
tightly connected the neighbors of a node are. High values indicate that
a node's neighbors are also likely to be connected, forming local
clusters of related symptoms.

\paragraph{Clustering Coefficient}\label{clustering-coefficient}

Clustering measures, such as Onnela, Zhang, and WS, are statistical
techniques used to identify groups or clusters within a dataset. These
methods help uncover hidden patterns and relationships between data
points. In the context of symptom networks, these measures can be used
to group similar symptoms together. Onnela's clustering coefficient is
more sensitive to the overall network structure, capturing subtle
clustering patterns (Constantini et al., 2015). Zhang and Horvath's
method focuses on identifying densely connected modules, while Watts and
Strogatz's transitivity emphasizes local clustering and the formation of
triangles(Constantini et al., 2015).

In our network, it is possible to see that a high clustering
coefficient, specially in WS and Zhang plots. All methods detect
consistent clusters for some items, especially in the SAS, QIDS, and
CESD scales, where items related to sleep, appetite/weight, and somatic
issues tend to vary more than the other items. We have already seen that
these type of items tend to be either peripheral (low strength and
closenness) or bridges the periphery (betweenness).

\begin{Shaded}
\begin{Highlighting}[]
\CommentTok{\# Creating a copy of the symptom similarity matrix}
\NormalTok{cl\_symp }\OtherTok{\textless{}{-}}\NormalTok{ symp\_spear}

\CommentTok{\# Making the matrix symmetric by copying the upper triangle to the lower triangle}
\NormalTok{cl\_symp[}\FunctionTok{upper.tri}\NormalTok{(cl\_symp)] }\OtherTok{\textless{}{-}} \FunctionTok{t}\NormalTok{(cl\_symp)[}\FunctionTok{upper.tri}\NormalTok{(cl\_symp)]}

\CommentTok{\# Creating a clustering plot based on the symmetric similarity matrix}
\FunctionTok{clusteringPlot}\NormalTok{(cl\_symp, }\AttributeTok{signed =}\NormalTok{ T, }\AttributeTok{decreasing =}\NormalTok{ T)}
\end{Highlighting}
\end{Shaded}

\includegraphics{network_analysis_files/figure-latex/unnamed-chunk-23-1.pdf}

\subsubsection{Density and Small-Worldness (Transitivity, APL - Average
Path Length, Small-World
Index):}\label{density-and-small-worldness-transitivity-apl---average-path-length-small-world-index}

Network density refers to the proportion of possible edges that are
actually present. Small-worldness combines transitivity (the likelihood
that two connected nodes share a neighbor) and average path length (the
average number of steps between two nodes) to determine if the network
has a small-world structure, where symptoms are highly clustered but
still have short paths between them.

A transitivity of 0.46 means that around 46\% of all possible
connections between nodes are present. This indicates a relatively
well-connected network but not saturated. The value of observed
transitivity (1) compared to a random network (0.99) suggests a strong
clustering structure. The APL 1 means that most nodes are directly
connected. It is slightly shorter than the APL in a random network
(1.38). The nodes are directly connected, which is slightly shorter than
the APL in a random network (1.38). ⁤⁤The small-world index of 1.39
indicates that the network combines high clustering with short paths
between nodes and this makes efficient connections within clusters
possible.

\paragraph{Density and smallworldness}\label{density-and-smallworldness}

\begin{Shaded}
\begin{Highlighting}[]
\CommentTok{\#Density: }
\FunctionTok{cat}\NormalTok{(}\StringTok{"Density: "}\NormalTok{, }\FunctionTok{mean}\NormalTok{(symp\_spear[}\FunctionTok{lower.tri}\NormalTok{(symp\_spear)]), }\StringTok{"}\SpecialCharTok{\textbackslash{}n}\StringTok{"}\NormalTok{)}
\end{Highlighting}
\end{Shaded}

\begin{verbatim}
## Density:  0.4647429
\end{verbatim}

\begin{Shaded}
\begin{Highlighting}[]
\CommentTok{\#APL, transitivity and small{-}worldness: }
\FunctionTok{smallworldIndex}\NormalTok{(}\FunctionTok{qgraph}\NormalTok{(symp\_spear))}
\end{Highlighting}
\end{Shaded}

\includegraphics{network_analysis_files/figure-latex/unnamed-chunk-24-1.pdf}

\begin{verbatim}
## $transitivity
## [1] 1
## 
## $transitivity_random
## [1] 0.9918699
## 
## $APL
## [1] 1
## 
## $APL_random
## [1] 1.381547
## 
## $index
## [1] 1.392871
\end{verbatim}

\subsection{Groups of symptoms based on core symptoms from
DSM-V-TR}\label{groups-of-symptoms-based-on-core-symptoms-from-dsm-v-tr}

\href{12\%20core\%20DSM-V-TR\%20symptoms\%20present\%20in\%20the\%20scales}{\includegraphics{images/venn2.PNG}}

descripción de los síntomas y su atribución

quizás añadir qgraph con grupos separados por síntomas

\subsection{Louvain Communities}\label{louvain-communities}

Explicar el concepto de louvain Explicar el concepto de modularidad

\subsubsection{Creating the absolute values adjacency
matrix}\label{creating-the-absolute-values-adjacency-matrix}

\begin{Shaded}
\begin{Highlighting}[]
\CommentTok{\# Converting the correlation matrix into an adjacency matrix}
\NormalTok{adj\_matrix }\OtherTok{\textless{}{-}} \FunctionTok{as.matrix}\NormalTok{(symp\_spear)}

\CommentTok{\# threshold to binarize the matrix (only keep strong correlations)}
\CommentTok{\#adj\_matrix[abs(adj\_matrix) \textless{} 0.3] \textless{}{-} 0  \# Set a threshold of 0.3 (or adjust as needed)}

\CommentTok{\# attribute abs values to negative edges}
\NormalTok{adj\_matrix[adj\_matrix }\SpecialCharTok{\textless{}} \DecValTok{0}\NormalTok{] }\OtherTok{\textless{}{-}} \FunctionTok{abs}\NormalTok{(adj\_matrix[adj\_matrix }\SpecialCharTok{\textless{}} \DecValTok{0}\NormalTok{])}

\CommentTok{\# Converting the adjacency matrix to an igraph object}
\NormalTok{graph\_symp }\OtherTok{\textless{}{-}} \FunctionTok{graph\_from\_adjacency\_matrix}\NormalTok{(adj\_matrix, }\AttributeTok{mode =} \StringTok{"undirected"}\NormalTok{, }\AttributeTok{weighted =} \ConstantTok{TRUE}\NormalTok{, }\AttributeTok{diag =} \ConstantTok{FALSE}\NormalTok{)}
\end{Highlighting}
\end{Shaded}

\subsubsection{Refining the resolution}\label{refining-the-resolution}

Changing resolution levels in order to find the most fine-grained set of
communities. Counting the number of symptoms represented in all scales
we arrived at 15 core symptoms related or not to DSM-V-TR. Using
resolutions between 1.05 to 1.1 creates communities from 8 to 17
communities. Therefore, the idea of employing this range comes from the
estimation of 15 main symptoms that we classified earlier. ** ESCRIBIR
SÍNTOMAS / PROFUNDIZAR

\begin{Shaded}
\begin{Highlighting}[]
\CommentTok{\# Running the Louvain algorithm 41 times to detect communities with a slowly growing modularity}
\NormalTok{decseq }\OtherTok{\textless{}{-}} \FunctionTok{seq}\NormalTok{(}\AttributeTok{from =} \FloatTok{1.05}\NormalTok{, }\AttributeTok{to =} \FloatTok{1.1}\NormalTok{, }\AttributeTok{by =} \FloatTok{0.00025}\NormalTok{)}
\NormalTok{templist }\OtherTok{\textless{}{-}} \FunctionTok{c}\NormalTok{()}

\CommentTok{\# For loop to get all results as new columns}
\ControlFlowTok{for}\NormalTok{(i }\ControlFlowTok{in}\NormalTok{ decseq)\{}
\NormalTok{  louvain\_commu\_temp }\OtherTok{\textless{}{-}} \FunctionTok{cluster\_louvain}\NormalTok{(graph\_symp, }\AttributeTok{resolution =}\NormalTok{ i)}
\NormalTok{  comm\_df }\OtherTok{\textless{}{-}} \FunctionTok{as.data.frame}\NormalTok{(}\FunctionTok{membership}\NormalTok{(louvain\_commu\_temp))}
  
\NormalTok{  templist[[}\FunctionTok{as.character}\NormalTok{(i)]] }\OtherTok{\textless{}{-}}\NormalTok{ comm\_df\}}

\NormalTok{louvain\_rep }\OtherTok{\textless{}{-}} \FunctionTok{as.data.frame}\NormalTok{(templist)}
\NormalTok{tlouvain }\OtherTok{\textless{}{-}} \FunctionTok{t}\NormalTok{(louvain\_rep)}


\CommentTok{\#write.xlsx(as.data.frame(tlouvain), "louvain.xlsx", rowNames = TRUE)}
\end{Highlighting}
\end{Shaded}

\subsubsection{Choosing the best
communities}\label{choosing-the-best-communities}

Find which result has the best correlation from all communities

\begin{Shaded}
\begin{Highlighting}[]
\CommentTok{\# Obtaining correlations}
\NormalTok{correlation\_matrix }\OtherTok{\textless{}{-}} \FunctionTok{cor}\NormalTok{(louvain\_rep, }\AttributeTok{method =} \StringTok{"pearson"}\NormalTok{)}
\NormalTok{clou }\OtherTok{\textless{}{-}} \FunctionTok{as.data.frame}\NormalTok{(correlation\_matrix)}

\CommentTok{\#}
\CommentTok{\#colnames(clou)}
\CommentTok{\# summing each column {-} summing is enough because all correlations are positive and have self cor values}
\NormalTok{l }\OtherTok{\textless{}{-}} \FunctionTok{list}\NormalTok{()}
\ControlFlowTok{for}\NormalTok{(i }\ControlFlowTok{in} \FunctionTok{colnames}\NormalTok{(clou))\{}
  \FunctionTok{cat}\NormalTok{(i, }\StringTok{" = "}\NormalTok{,}\FunctionTok{sum}\NormalTok{(clou[[i]]), }\StringTok{"}\SpecialCharTok{\textbackslash{}n}\StringTok{"}\NormalTok{)}
\NormalTok{  l }\OtherTok{\textless{}{-}} \FunctionTok{append}\NormalTok{(l, }\FunctionTok{sum}\NormalTok{(clou[[i]]))}
\NormalTok{\}}
\end{Highlighting}
\end{Shaded}

\begin{verbatim}
## x  =  150.0512 
## x.1  =  86.91847 
## x.2  =  145.0386 
## x.3  =  148.1955 
## x.4  =  144.4036 
## x.5  =  148.1955 
## x.6  =  166.0505 
## x.7  =  144.4036 
## x.8  =  144.4036 
## x.9  =  145.0386 
## x.10  =  151.2141 
## x.11  =  106.0144 
## x.12  =  98.21327 
## x.13  =  106.0144 
## x.14  =  141.7053 
## x.15  =  139.9899 
## x.16  =  125.9182 
## x.17  =  144.4036 
## x.18  =  145.0386 
## x.19  =  146.3934 
## x.20  =  137.4432 
## x.21  =  153.06 
## x.22  =  141.3561 
## x.23  =  159.9188 
## x.24  =  164.9384 
## x.25  =  102.0587 
## x.26  =  106.0144 
## x.27  =  98.00002 
## x.28  =  166.5386 
## x.29  =  159.9952 
## x.30  =  153.06 
## x.31  =  141.7234 
## x.32  =  148.3663 
## x.33  =  146.3934 
## x.34  =  139.1027 
## x.35  =  147.8963 
## x.36  =  146.3934 
## x.37  =  151.2141 
## x.38  =  168.218 
## x.39  =  165.8145 
## x.40  =  148.4588 
## x.41  =  142.8205 
## x.42  =  138.6498 
## x.43  =  142.8205 
## x.44  =  165.8145 
## x.45  =  143.8279 
## x.46  =  142.8205 
## x.47  =  168.218 
## x.48  =  166.8474 
## x.49  =  151.2141 
## x.50  =  142.8205 
## x.51  =  145.8261 
## x.52  =  160.1606 
## x.53  =  166.1656 
## x.54  =  142.4329 
## x.55  =  169.1198 
## x.56  =  165.8751 
## x.57  =  169.1246 
## x.58  =  167.9812 
## x.59  =  163.7674 
## x.60  =  164.9352 
## x.61  =  156.6877 
## x.62  =  148.5018 
## x.63  =  163.21 
## x.64  =  148.5018 
## x.65  =  144.6695 
## x.66  =  162.4698 
## x.67  =  150.0976 
## x.68  =  150.3435 
## x.69  =  132.1573 
## x.70  =  163.21 
## x.71  =  143.8915 
## x.72  =  162.36 
## x.73  =  135.2534 
## x.74  =  163.555 
## x.75  =  164.1931 
## x.76  =  165.5293 
## x.77  =  165.5293 
## x.78  =  46.19869 
## x.79  =  163.9511 
## x.80  =  168.9856 
## x.81  =  171.3497 
## x.82  =  166.4445 
## x.83  =  131.5738 
## x.84  =  168.7345 
## x.85  =  141.7079 
## x.86  =  168.3634 
## x.87  =  167.5508 
## x.88  =  167.1203 
## x.89  =  168.3195 
## x.90  =  166.7079 
## x.91  =  167.6884 
## x.92  =  168.3634 
## x.93  =  167.6884 
## x.94  =  168.7767 
## x.95  =  161.9512 
## x.96  =  170.2783 
## x.97  =  167.6884 
## x.98  =  167.253 
## x.99  =  170.9212 
## x.100  =  170.2783 
## x.101  =  162.9905 
## x.102  =  167.5244 
## x.103  =  165.7521 
## x.104  =  168.0643 
## x.105  =  169.6433 
## x.106  =  164.2818 
## x.107  =  172.2502 
## x.108  =  165.1378 
## x.109  =  171.0213 
## x.110  =  168.3792 
## x.111  =  168.9202 
## x.112  =  170.1779 
## x.113  =  170.1844 
## x.114  =  167.0424 
## x.115  =  164.7249 
## x.116  =  169.3205 
## x.117  =  142.0456 
## x.118  =  165.1007 
## x.119  =  169.5734 
## x.120  =  170.1596 
## x.121  =  170.4229 
## x.122  =  166.3972 
## x.123  =  162.2151 
## x.124  =  136.8291 
## x.125  =  168.403 
## x.126  =  172.4335 
## x.127  =  171.6411 
## x.128  =  162.7232 
## x.129  =  170.7825 
## x.130  =  64.92481 
## x.131  =  171.6411 
## x.132  =  172.3611 
## x.133  =  171.1987 
## x.134  =  170.9784 
## x.135  =  156.7919 
## x.136  =  169.7458 
## x.137  =  172.6135 
## x.138  =  172.9804 
## x.139  =  173.2353 
## x.140  =  171.9489 
## x.141  =  162.807 
## x.142  =  149.874 
## x.143  =  171.998 
## x.144  =  171.0238 
## x.145  =  163.4276 
## x.146  =  170.9068 
## x.147  =  169.7866 
## x.148  =  88.51036 
## x.149  =  157.529 
## x.150  =  162.954 
## x.151  =  171.448 
## x.152  =  164.7388 
## x.153  =  169.8647 
## x.154  =  165.1064 
## x.155  =  172.8904 
## x.156  =  171.6401 
## x.157  =  171.5207 
## x.158  =  170.284 
## x.159  =  162.4284 
## x.160  =  171.1102 
## x.161  =  152.1521 
## x.162  =  171.2214 
## x.163  =  155.6202 
## x.164  =  169.3939 
## x.165  =  118.1083 
## x.166  =  168.2039 
## x.167  =  166.5189 
## x.168  =  172.1858 
## x.169  =  170.284 
## x.170  =  157.0715 
## x.171  =  152.1581 
## x.172  =  166.8646 
## x.173  =  173.0115 
## x.174  =  167.7343 
## x.175  =  171.1017 
## x.176  =  170.9206 
## x.177  =  172.4764 
## x.178  =  136.9115 
## x.179  =  173.0697 
## x.180  =  159.0943 
## x.181  =  163.453 
## x.182  =  157.4526 
## x.183  =  166.7138 
## x.184  =  171.3073 
## x.185  =  157.6184 
## x.186  =  167.3521 
## x.187  =  157.5549 
## x.188  =  157.5549 
## x.189  =  156.776 
## x.190  =  156.8088 
## x.191  =  170.3032 
## x.192  =  167.6603 
## x.193  =  173.2236 
## x.194  =  170.3484 
## x.195  =  162.0409 
## x.196  =  169.9611 
## x.197  =  173.8528 
## x.198  =  165.7221 
## x.199  =  163.746 
## x.200  =  160.7241
\end{verbatim}

\begin{Shaded}
\begin{Highlighting}[]
\CommentTok{\#value with the strongest correlations from all}
\NormalTok{clou[}\FunctionTok{which.max}\NormalTok{(}\FunctionTok{unlist}\NormalTok{(l))] }
\end{Highlighting}
\end{Shaded}

\begin{verbatim}
##           x.197
## x     0.7829956
## x.1   0.5037434
## x.2   0.7499704
## x.3   0.7730168
## x.4   0.7498843
## x.5   0.7730168
## x.6   0.9304422
## x.7   0.7498843
## x.8   0.7498843
## x.9   0.7499704
## x.10  0.7911360
## x.11  0.5464325
## x.12  0.4957089
## x.13  0.5464325
## x.14  0.7235154
## x.15  0.7174582
## x.16  0.6953857
## x.17  0.7498843
## x.18  0.7499704
## x.19  0.7680693
## x.20  0.7537533
## x.21  0.8053980
## x.22  0.7245999
## x.23  0.8712440
## x.24  0.9022109
## x.25  0.5201542
## x.26  0.5464325
## x.27  0.5024572
## x.28  0.9313911
## x.29  0.8792246
## x.30  0.8053980
## x.31  0.7226016
## x.32  0.7756145
## x.33  0.7680693
## x.34  0.7089456
## x.35  0.7783104
## x.36  0.7680693
## x.37  0.7911360
## x.38  0.9295090
## x.39  0.9208136
## x.40  0.7836248
## x.41  0.7360802
## x.42  0.7671673
## x.43  0.7360802
## x.44  0.9208136
## x.45  0.7508455
## x.46  0.7360802
## x.47  0.9295090
## x.48  0.9176257
## x.49  0.7911360
## x.50  0.7360802
## x.51  0.7605759
## x.52  0.8821910
## x.53  0.9309929
## x.54  0.7392323
## x.55  0.9498472
## x.56  0.9175764
## x.57  0.9592605
## x.58  0.9106184
## x.59  0.9222725
## x.60  0.9304522
## x.61  0.8443233
## x.62  0.7820952
## x.63  0.8915262
## x.64  0.7820952
## x.65  0.7695370
## x.66  0.8993878
## x.67  0.7929125
## x.68  0.8028103
## x.69  0.7450194
## x.70  0.8915262
## x.71  0.7551372
## x.72  0.8874652
## x.73  0.7571388
## x.74  0.9232878
## x.75  0.9279423
## x.76  0.9235413
## x.77  0.9235413
## x.78  0.3002378
## x.79  0.9143615
## x.80  0.9458373
## x.81  0.9675864
## x.82  0.9179136
## x.83  0.7415721
## x.84  0.9512298
## x.85  0.7837018
## x.86  0.9506618
## x.87  0.9301278
## x.88  0.9188082
## x.89  0.9432981
## x.90  0.9195775
## x.91  0.9333751
## x.92  0.9506618
## x.93  0.9333751
## x.94  0.9665235
## x.95  0.8748466
## x.96  0.9572875
## x.97  0.9333751
## x.98  0.9356783
## x.99  0.9664379
## x.100 0.9572875
## x.101 0.8860249
## x.102 0.9351468
## x.103 0.9070355
## x.104 0.9306148
## x.105 0.9317327
## x.106 0.9101521
## x.107 0.9748253
## x.108 0.9338949
## x.109 0.9610248
## x.110 0.9217427
## x.111 0.9396234
## x.112 0.9701596
## x.113 0.9469250
## x.114 0.9160914
## x.115 0.9010476
## x.116 0.9378564
## x.117 0.8021420
## x.118 0.9201712
## x.119 0.9366672
## x.120 0.9352058
## x.121 0.9458122
## x.122 0.9266496
## x.123 0.8915918
## x.124 0.7628464
## x.125 0.9408622
## x.126 0.9763107
## x.127 0.9640545
## x.128 0.9027981
## x.129 0.9559378
## x.130 0.3920193
## x.131 0.9640545
## x.132 0.9770469
## x.133 0.9695968
## x.134 0.9599702
## x.135 0.8727690
## x.136 0.9429876
## x.137 0.9759135
## x.138 0.9843482
## x.139 0.9866087
## x.140 0.9846479
## x.141 0.9034364
## x.142 0.8344740
## x.143 0.9756615
## x.144 0.9514003
## x.145 0.9144407
## x.146 0.9511007
## x.147 0.9592542
## x.148 0.5251412
## x.149 0.8718467
## x.150 0.9145240
## x.151 0.9587377
## x.152 0.9073647
## x.153 0.9605651
## x.154 0.9086929
## x.155 0.9845725
## x.156 0.9634679
## x.157 0.9754583
## x.158 0.9635223
## x.159 0.8937592
## x.160 0.9521007
## x.161 0.8436435
## x.162 0.9516278
## x.163 0.8599048
## x.164 0.9557214
## x.165 0.6597321
## x.166 0.9288137
## x.167 0.9200436
## x.168 0.9856871
## x.169 0.9635223
## x.170 0.8694018
## x.171 0.8323557
## x.172 0.9265951
## x.173 0.9849569
## x.174 0.9349802
## x.175 0.9541966
## x.176 0.9492098
## x.177 0.9811801
## x.178 0.7843518
## x.179 0.9854122
## x.180 0.8840017
## x.181 0.8945686
## x.182 0.8733809
## x.183 0.9315037
## x.184 0.9538465
## x.185 0.8456625
## x.186 0.9347515
## x.187 0.8741251
## x.188 0.8741251
## x.189 0.8719571
## x.190 0.8762889
## x.191 0.9764572
## x.192 0.9396639
## x.193 0.9895386
## x.194 0.9743919
## x.195 0.8832310
## x.196 0.9513738
## x.197 1.0000000
## x.198 0.9286463
## x.199 0.9105485
## x.200 0.8657738
\end{verbatim}

\begin{Shaded}
\begin{Highlighting}[]
\NormalTok{modularity }\OtherTok{\textless{}{-}}\NormalTok{ decseq[}\DecValTok{145}\NormalTok{]}
\NormalTok{found\_solution }\OtherTok{\textless{}{-}}\NormalTok{ louvain\_rep[}\DecValTok{145}\NormalTok{]}
\end{Highlighting}
\end{Shaded}

Having found the fittest modularity, we could iterate it in order to
find a consensus.

\begin{Shaded}
\begin{Highlighting}[]
\CommentTok{\#comm\_consensus \textless{}{-} community.consensus(graph\_symp, resolution = modularity)}
\CommentTok{\#cons \textless{}{-} as.data.frame(comm\_consensus)}
\end{Highlighting}
\end{Shaded}

\subsection{Walktrap Communities}\label{walktrap-communities}

Explicar el concepto de walktrap

\begin{Shaded}
\begin{Highlighting}[]
\NormalTok{wt }\OtherTok{\textless{}{-}} \FunctionTok{cluster\_walktrap}\NormalTok{(graph\_symp)}
\NormalTok{dfwt }\OtherTok{\textless{}{-}} \FunctionTok{as.data.frame}\NormalTok{(wt}\SpecialCharTok{$}\NormalTok{membership)}
\end{Highlighting}
\end{Shaded}

\begin{Shaded}
\begin{Highlighting}[]
\NormalTok{df\_comm }\OtherTok{\textless{}{-}} \FunctionTok{read.csv}\NormalTok{(}\StringTok{\textquotesingle{}../data/communities.csv\textquotesingle{}}\NormalTok{)}
\NormalTok{df\_comm}
\end{Highlighting}
\end{Shaded}

\begin{verbatim}
##       n
## 1    10
## 2    21
## 3    53
## 4    95
## 5    70
## 6    85
## 7    83
## 8     1
## 9    18
## 10   63
## 11   65
## 12   25
## 13   64
## 14   66
## 15    3
## 16   20
## 17   27
## 18   67
## 19   45
## 20   82
## 21   87
## 22   36
## 23    2
## 24    4
## 25    6
## 26   22
## 27   44
## 28   62
## 29    8
## 30   37
## 31   39
## 32   46
## 33   48
## 34   31
## 35   33
## 36   34
## 37   40
## 38   43
## 39   47
## 40   50
## 41   51
## 42   52
## 43   54
## 44   55
## 45   97
## 46   99
## 47  100
## 48  101
## 49   96
## 50   84
## 51   81
## 52   16
## 53   24
## 54    7
## 55   71
## 56   57
## 57  115
## 58   72
## 59   98
## 60   77
## 61   78
## 62   80
## 63   89
## 64   94
## 65  108
## 66  110
## 67  107
## 68  111
## 69   28
## 70   13
## 71   12
## 72  112
## 73   73
## 74   88
## 75   90
## 76   91
## 77   93
## 78  102
## 79  103
## 80  104
## 81  105
## 82  106
## 83  109
## 84   11
## 85   61
## 86    5
## 87   19
## 88   23
## 89   58
## 90   59
## 91   60
## 92  123
## 93   76
## 94  113
## 95  118
## 96  114
## 97  121
## 98   74
## 99   75
## 100  79
## 101  92
## 102 116
## 103 117
## 104 119
## 105 120
## 106 122
## 107  14
## 108  15
## 109  29
## 110  69
## 111  49
## 112   9
## 113  17
## 114  26
## 115  30
## 116  35
## 117  38
## 118  41
## 119  56
## 120  68
## 121  32
## 122  42
## 123  86
##                                                                                                                                                     question
## 1                                                                                                           Había perdido interés en sus actividades diarias
## 2                                                                                                     Tener poco interés o disfrutar poco haciendo las cosas
## 3                                                                                                                               Sentí que nada me interesaba
## 4                                                                                                               Todavía disfruto con lo que antes me gustaba
## 5                                                                                                                                            Interés general
## 6                                                                                                                  He sido incapaz de entusiasmarse por nada
## 7                                                                                                             He sentido que no había nada que me ilusionara
## 8                                                                                                                                              Tenía apetito
## 9                                                                                                                                 Perdió peso sin intentarlo
## 10                                                                                                                         Ha habido disminución del apetito
## 11                                                                                                                                 Ha habido pérdida de peso
## 12                                                                                                                      Tener poco apetito o comer en exceso
## 13                                                                                                                             Ha habido aumento del apetito
## 14                                                                                                                                 Ha habido aumento de peso
## 15                                                                                         Tenía dificultad para mantener su mente en lo que estaba haciendo
## 16                                                                                                         Era difícil concentrarse en las cosas importantes
## 17                                                            Tener dificultades para concentrarse en cosas tales como leer el periódico o ver la televisión
## 18                                                                                                                          Concentración/Toma de decisiones
## 19                                                                                                                      Tuve problemas para tomar decisiones
## 20                                                                                                 Me ha resultado difícil tener iniciativa para hacer cosas
## 21                                                                                                            He sentido que la vida no tenía ningún sentido
## 22                                                                                                                                           Me sentí triste
## 23                                                                                                                             No podía quitarse la tristeza
## 24                                                                                                                                     Se sentía deprimido/a
## 25                                                                                                                                          Se sintió triste
## 26                                                                                                        Sentirse desanimado/a, deprimido/a o sin esperanza
## 27                                                                                                                                      Me sentí deprimido/a
## 28                                                                                                                                           Sentirse triste
## 29                                                                                                                                       Nada le hacía feliz
## 30                                                                                                              Sentí que quería darme por vencido/a en todo
## 31                                                                                                                                      Me sentí fracasado/a
## 32                                                                                                                   Me sentí desanimado/a acerca del futuro
## 33                                                                                                                                    Me sentí descontento/a
## 34                                                                                                                              Sentí que nada me ilusionaba
## 35                                                                                                                                 Me distancié de los demás
## 36                                                                                                                            Sentí que nada me podía animar
## 37                                                                                      Me costó trabajo sentir una relación estrecha (cercana) con la gente
## 38                                                                                                                                           Me sentí solo/a
## 39                                                                                            Descubrí que había cosas en mi vida que eran demasiado para mí
## 40                                                                                                                                 Me sentí desesperanzado/a
## 41                                                                                                                       Sentí que la gente no me hacía caso
## 42                                                                                                                   Me sentí angustiado/a sin ningún motivo
## 43                                                                                                  Me sentí pesimista (que veía las cosas de modo negativo)
## 44                                                                                                                            Sentí que mi vida estaba vacía
## 45                                                                                                                                          Me siento alegre
## 46                                                                                                                 He perdido interés en mi aspecto personal
## 47                                                                                                                    Me siento optimista respecto al futuro
## 48                                                                                      Me divierto con un buen libro, la radio, o un programa de televisión
## 49                                                                                                         Puedo reírme y ver el lado divertido de las cosas
## 50                                                                                                                       Me he sentido desanimado/a y triste
## 51                                                                                                              No he podido sentir ninguna emoción positiva
## 52                                                                                                                        Se sentía cansado/a todo el tiempo
## 53                                                                                                                   Sentirse cansado/a o tener poca energía
## 54                                                                                                                              No tenía ganas de hacer nada
## 55                                                                                                                                          Nivel de energía
## 56                                                                                                           Me sentí completamente agotado/a emocionalmente
## 57                                                                                                                     Me siento débil y me canso fácilmente
## 58                                                                                                                             Sentimiento de mayor lentitud
## 59                                                                                                         Me siento como si cada día estuviera más lento/a.
## 60                                                        He estado preocupado/a por situaciones en las que pudiera ser presa del pánico y hacer el ridículo
## 61                                                                                                                 He sentido que estaba al borde del pánico
## 62                                                                                                            Me he sentido asustado/a sin una razón de peso
## 63                                                                                    Tengo una sensación de miedo, como si algo horrible me fuera a suceder
## 64                                                                                                              Me asaltan sentimientos repentinos de pánico
## 65                                                                                                    Ha tenido miedo, como si fuera a suceder algo terrible
## 66                                                                                                                              Me da miedo sin razón alguna
## 67                                                                                                                   Se ha irritado o enfadado con facilidad
## 68                                                                                                                       Me enojo fácilmente o siento pánico
## 69  Moverse o hablar tan lentamente que los demás lo han notado. O bien al contrario, estar tan agitado/a o inquieto/a que se mueve mucho más de lo habitual
## 70                                                                                                                                       Se sentía agitado/a
## 71                                                                                                                             Sentía que se movía muy lento
## 72                                                                                                                          Siento que me estoy desmoronando
## 73                                                                                                                                    Sensación de inquietud
## 74                                                                                                                            Me siento tenso/a o nervioso/a
## 75                                                                                                                    Tengo mi mente llena de preocupaciones
## 76                                                                                              Puedo estar sentado/a confortablemente y sentirme relajado/a
## 77                                                                                                 Me siento inquieto/a, como si no pudiera parar de moverme
## 78                                                                                                      Se ha sentido nervioso/a, ansioso/a o muy alterado/a
## 79                                                                                                                         No ha podido dejar de preocuparse
## 80                                                                                                       Se ha preocupado excesivamente por diferentes cosas
## 81                                                                                                                       Ha tenido dificultad para relajarse
## 82                                                                                        Se ha sentido tan intranquilo/a que no ha podido quedarse quieto/a
## 83                                                                                                     Me siento más nervioso/a y ansioso/a que de costumbre
## 84                                                                                                                                 Dormía más de lo habitual
## 85                                                                                                                                          Dormir demasiado
## 86                                                                                                                                     Su sueño era inquieto
## 87                                                                                                                           Le costaba mucho trabajo dormir
## 88                                                                                      Tener problemas para dormir o mantenerse dormido, o dormir demasiado
## 89                                                                                                                           Problemas a la hora de dormirse
## 90                                                                                                                                   Sueño nocturno alterado
## 91                                                                                                                            Despertarse demasiado temprano
## 92                                                                                                                                          Tengo pesadillas
## 93                                                                                                                 He tenido temblores (p.ej., en las manos)
## 94                                                                                                                             Mis brazos y piernas tiemblan
## 95                                                                                                                    Tengo desmayos o sensación de desmayos
## 96                                                                                              Me molestan los dolores de cabeza, cuello y dolor de espalda
## 97                                                                                                                 Tengo que vaciar mi vejiga con frecuencia
## 98                                                                                                                                    He notado la boca seca
## 99                          He tenido dificultades para respirar (p.ej., respiración excesivamente rápida, falta de aliento sin haber hecho esfuerzo físico)
## 100                                             He notado mi corazón sin hacer esfuerzo físico (p.ej., aumento del ritmo cardiaco, ausencia de algún latido)
## 101                                                                                    Tengo una sensación extraña, como si tuviera mariposas en el estómago
## 102                                                                                                                     Puedo sentir mi corazón latir rápido
## 103                                                                                                                                   Me molestan los mareos
## 104                                                                               Tengo sensación de entumecimiento y hormigueo en mis dedos de manos y pies
## 105                                                                                                     Me molestan los dolores de estómago o la indigestión
## 106                                                                                                           Se me calienta la cara y se me pone enrojecida
## 107                                                                                                                          Sentía deseos de estar muerto/a
## 108                                                                                                                                      Quería hacerse daño
## 109                                                             Tener pensamientos sobre estar muerto(a) o sobre hacerse daño a sí mismo(a) de alguna manera
## 110                                                                                                               Pensamientos sobre la muerte o el suicidio
## 111                                                                                                              Sentí que no tenía ninguna razón para vivir
## 112                                                                                                                          Sentía que era una mala persona
## 113                                                                                                                        Estaba a disgusto consigo mismo/a
## 114                                  Sentirse mal consigo mismo/a o sentirse fracasado/a, o pensar que se ha decepcionado a sí mismo/a o a los que le rodean
## 115                                                                                                                                  Sentí que no valía nada
## 116                                                                                                             Sentí que no valía tanto como otras personas
## 117                                                                                                                Sentí que tenía la culpa de lo que pasaba
## 118                                                                                                                    Me sentí decepcionado/a de mí mismo/a
## 119                                                                                                                                        Me sentí culpable
## 120                                                                                                                                    Opinión de sí mismo/a
## 121                                                                                             Me sentí indefenso/a (que no podía hacer nada para ayudarme)
## 122                                                                                                                            Sentí que nadie me necesitaba
## 123                                                                                                               He sentido que no valía mucho como persona
##         items     core_dsm_symptoms louvain_comm_consensus_r1.086
## 1     cesd_10             anhedonia                             3
## 2      phq_01             anhedonia                             3
## 3   promis_24             anhedonia                             7
## 4     hads_02             anhedonia                             8
## 5     qids_13             anhedonia                             8
## 6     dass_16             anhedonia                            11
## 7     dass_10             anhedonia                            11
## 8     cesd_01        apetite/weight                             1
## 9     cesd_18        apetite/weight                             1
## 10    qids_06        apetite/weight                             1
## 11    qids_08        apetite/weight                             1
## 12     phq_05        apetite/weight                             6
## 13    qids_07        apetite/weight                             6
## 14    qids_09        apetite/weight                             6
## 15    cesd_03         concentration                             3
## 16    cesd_20         concentration                             3
## 17     phq_07         concentration                             3
## 18    qids_10         concentration                             6
## 19  promis_16         concentration                             7
## 20    dass_05         concentration                            11
## 21    dass_21        depressed mood                             1
## 22  promis_07        depressed mood                             2
## 23    cesd_02        depressed mood                             2
## 24    cesd_04        depressed mood                             2
## 25    cesd_06        depressed mood                             2
## 26     phq_02        depressed mood                             2
## 27  promis_15        depressed mood                             2
## 28    qids_05        depressed mood                             2
## 29    cesd_08        depressed mood                             3
## 30  promis_08        depressed mood                             5
## 31  promis_10        depressed mood                             5
## 32  promis_17        depressed mood                             7
## 33  promis_19        depressed mood                             7
## 34  promis_02        depressed mood                             7
## 35  promis_04        depressed mood                             7
## 36  promis_05        depressed mood                             7
## 37  promis_11        depressed Mood                             7
## 38  promis_14        depressed mood                             7
## 39  promis_18        depressed mood                             7
## 40  promis_21        depressed mood                             7
## 41  promis_22        depressed mood                             7
## 42  promis_23        depressed mood                             7
## 43  promis_25        depressed mood                             7
## 44  promis_26        depressed mood                             7
## 45    hads_06        depressed mood                             8
## 46    hads_10        depressed mood                             8
## 47    hads_12        depressed mood                             8
## 48    hads_14        depressed mood                             8
## 49    hads_04        depressed mood                             8
## 50    dass_13        depressed mood                            11
## 51    dass_03        depressed mood                            11
## 52    cesd_16               fatigue                             3
## 53     phq_04               fatigue                             3
## 54    cesd_07               fatigue                             3
## 55    qids_14               fatigue                             6
## 56  promis_28               fatigue                             7
## 57     sas_07               fatigue                             9
## 58    qids_15               fatigue                             9
## 59    hads_08               fatigue                             9
## 60    dass_09         fear/keyed up                             9
## 61    dass_15         fear/keyed up                             9
## 62    dass_20         fear/keyed up                             9
## 63    hads_03         fear/keyed up                             9
## 64    hads_13         fear/keyed up                             9
## 65     gad_07         fear/keyed up                             9
## 66     sas_02         fear/keyed up                             9
## 67     gad_06          Irritability                            10
## 68     sas_03          Irritability                            10
## 69     phq_08           psychomotor                             1
## 70    cesd_13           psychomotor                             3
## 71    cesd_12           psychomotor                             3
## 72     sas_04 restlessness/keyed up                             9
## 73    qids_16 restlessness/keyed up                            10
## 74    hads_01 restlessness/keyed up                            10
## 75    hads_05 restlessness/keyed up                            10
## 76    hads_07 restlessness/keyed up                            10
## 77    hads_11 restlessness/keyed up                            10
## 78     gad_01 restlessness/keyed up                            10
## 79     gad_02 restlessness/keyed up                            10
## 80     gad_03 restlessness/keyed up                            10
## 81     gad_04 restlessness/keyed up                            10
## 82     gad_05 restlessness/keyed up                            10
## 83     sas_01 restlessness/keyed up                            10
## 84    cesd_11                 sleep                             1
## 85    qids_04                 sleep                             1
## 86    cesd_05                 sleep                             4
## 87    cesd_19                 sleep                             4
## 88     phq_03                 sleep                             4
## 89    qids_01                 sleep                             4
## 90    qids_02                 sleep                             4
## 91    qids_03                 sleep                             4
## 92     sas_15                 sleep                             9
## 93    dass_07      somatic symptoms                             9
## 94     sas_05      somatic symptoms                             9
## 95     sas_10      somatic symptoms                             9
## 96     sas_06      somatic symptoms                             9
## 97     sas_13      somatic symptoms                             9
## 98    dass_02      somatic symptoms                             9
## 99    dass_04      somatic symptoms                             9
## 100   dass_19      somatic symptoms                             9
## 101   hads_09      somatic symptoms                             9
## 102    sas_08      somatic symptoms                             9
## 103    sas_09      somatic symptoms                             9
## 104    sas_11      somatic symptoms                             9
## 105    sas_12      somatic symptoms                             9
## 106    sas_14      somatic symptoms                             9
## 107   cesd_14               suicide                             1
## 108   cesd_15               suicide                             1
## 109    phq_09               suicide                             1
## 110   qids_12               suicide                             1
## 111 promis_20               suicide                             1
## 112   cesd_09   worthlessness/guilt                             1
## 113   cesd_17   worthlessness/guilt                             5
## 114    phq_06   worthlessness/guilt                             5
## 115 promis_01   worthlessness/guilt                             5
## 116 promis_06   worthlessness/guilt                             5
## 117 promis_09   worthlessness/guilt                             5
## 118 promis_12   worthlessness/guilt                             5
## 119 promis_27   worthlessness/guilt                             5
## 120   qids_11   worthlessness/guilt                             6
## 121 promis_03   worthlessness/guilt                             7
## 122 promis_13   worthlessness/guilt                             7
## 123   dass_17   worthlessness/guilt                            11
##     walktrap_communities
## 1                      4
## 2                      4
## 3                      4
## 4                      2
## 5                      4
## 6                      3
## 7                      4
## 8                      1
## 9                      1
## 10                     1
## 11                     1
## 12                     2
## 13                     2
## 14                     2
## 15                     4
## 16                     4
## 17                     4
## 18                     2
## 19                     4
## 20                     2
## 21                     3
## 22                     2
## 23                     4
## 24                     4
## 25                     4
## 26                     4
## 27                     4
## 28                     4
## 29                     4
## 30                     4
## 31                     4
## 32                     2
## 33                     2
## 34                     4
## 35                     4
## 36                     4
## 37                     4
## 38                     4
## 39                     4
## 40                     4
## 41                     4
## 42                     4
## 43                     4
## 44                     4
## 45                     2
## 46                     2
## 47                     2
## 48                     3
## 49                     4
## 50                     2
## 51                     3
## 52                     2
## 53                     2
## 54                     4
## 55                     2
## 56                     2
## 57                     2
## 58                     3
## 59                     3
## 60                     3
## 61                     3
## 62                     3
## 63                     3
## 64                     3
## 65                     3
## 66                     3
## 67                     2
## 68                     2
## 69                     3
## 70                     2
## 71                     3
## 72                     3
## 73                     2
## 74                     2
## 75                     2
## 76                     2
## 77                     2
## 78                     2
## 79                     2
## 80                     2
## 81                     2
## 82                     2
## 83                     2
## 84                     1
## 85                     1
## 86                     2
## 87                     2
## 88                     2
## 89                     2
## 90                     2
## 91                     3
## 92                     3
## 93                     1
## 94                     1
## 95                     1
## 96                     2
## 97                     2
## 98                     3
## 99                     3
## 100                    3
## 101                    3
## 102                    3
## 103                    3
## 104                    3
## 105                    3
## 106                    3
## 107                    1
## 108                    1
## 109                    1
## 110                    1
## 111                    3
## 112                    1
## 113                    4
## 114                    4
## 115                    4
## 116                    4
## 117                    4
## 118                    4
## 119                    4
## 120                    4
## 121                    4
## 122                    4
## 123                    4
\end{verbatim}

\begin{Shaded}
\begin{Highlighting}[]
\NormalTok{wt1 }\OtherTok{\textless{}{-}}\NormalTok{ df\_comm[df\_comm}\SpecialCharTok{$}\NormalTok{walktrap\_communities }\SpecialCharTok{==} \DecValTok{1}\NormalTok{, ]}\SpecialCharTok{$}\NormalTok{n}
\NormalTok{wt2 }\OtherTok{\textless{}{-}}\NormalTok{ df\_comm[df\_comm}\SpecialCharTok{$}\NormalTok{walktrap\_communities }\SpecialCharTok{==} \DecValTok{2}\NormalTok{, ]}\SpecialCharTok{$}\NormalTok{n}
\NormalTok{wt3 }\OtherTok{\textless{}{-}}\NormalTok{ df\_comm[df\_comm}\SpecialCharTok{$}\NormalTok{walktrap\_communities }\SpecialCharTok{==} \DecValTok{3}\NormalTok{, ]}\SpecialCharTok{$}\NormalTok{n}
\NormalTok{wt4 }\OtherTok{\textless{}{-}}\NormalTok{ df\_comm[df\_comm}\SpecialCharTok{$}\NormalTok{walktrap\_communities }\SpecialCharTok{==} \DecValTok{4}\NormalTok{, ]}\SpecialCharTok{$}\NormalTok{n}



\NormalTok{color\_nodes\_symp\_df }\OtherTok{\textless{}{-}} \FunctionTok{c}\NormalTok{(}\StringTok{"darkblue"}\NormalTok{, }\StringTok{"limegreen"}\NormalTok{, }\StringTok{"darkorange"}\NormalTok{, }\StringTok{"purple"}\NormalTok{) }\CommentTok{\#, "grey", "yellow", "black", "black", "magenta", "brown", "darkorchid4" ,"white", "black" )}


\FunctionTok{qgraph}\NormalTok{(symp\_spear, }\AttributeTok{layout=}\StringTok{"spring"}\NormalTok{, }\AttributeTok{labels =} \FunctionTok{colnames}\NormalTok{(readable\_anx\_dp), }\AttributeTok{sampleSize =} \FunctionTok{nrow}\NormalTok{(df\_anx\_dp\_adj), }\AttributeTok{groups =} \FunctionTok{list}\NormalTok{(}\StringTok{\textquotesingle{}C1 {-} Sui/Sleep/Soma/Appe\textquotesingle{}} \OtherTok{=}\NormalTok{ wt1, }\StringTok{\textquotesingle{}C2 {-} Sleep/Fatig/Rest/Dep\textquotesingle{}} \OtherTok{=}\NormalTok{ wt2, }\StringTok{\textquotesingle{}C3 {-} Soma/Fear/Moto/Rest\textquotesingle{}} \OtherTok{=}\NormalTok{ wt3, }\StringTok{\textquotesingle{}C4 {-} Anhe/Dep/Wort/Conc\textquotesingle{}} \OtherTok{=}\NormalTok{ wt4), }\AttributeTok{color =}\NormalTok{ color\_nodes\_symp\_df, }\AttributeTok{theme=}\StringTok{"TeamFortress"}\NormalTok{, }\AttributeTok{label.prop=}\FloatTok{0.8}\NormalTok{, }\AttributeTok{vsize=}\DecValTok{2}\NormalTok{, }\AttributeTok{label.cex=}\DecValTok{1}\NormalTok{)}
\end{Highlighting}
\end{Shaded}

\includegraphics{network_analysis_files/figure-latex/unnamed-chunk-31-1.pdf}

\begin{Shaded}
\begin{Highlighting}[]
\NormalTok{lv1 }\OtherTok{\textless{}{-}}\NormalTok{ df\_comm[df\_comm}\SpecialCharTok{$}\NormalTok{louvain\_comm\_consensus\_r1}\FloatTok{.086} \SpecialCharTok{==} \DecValTok{1}\NormalTok{, ]}\SpecialCharTok{$}\NormalTok{n}
\NormalTok{lv2 }\OtherTok{\textless{}{-}}\NormalTok{ df\_comm[df\_comm}\SpecialCharTok{$}\NormalTok{louvain\_comm\_consensus\_r1}\FloatTok{.086} \SpecialCharTok{==} \DecValTok{2}\NormalTok{, ]}\SpecialCharTok{$}\NormalTok{n}
\NormalTok{lv3 }\OtherTok{\textless{}{-}}\NormalTok{ df\_comm[df\_comm}\SpecialCharTok{$}\NormalTok{louvain\_comm\_consensus\_r1}\FloatTok{.086} \SpecialCharTok{==} \DecValTok{3}\NormalTok{, ]}\SpecialCharTok{$}\NormalTok{n}
\NormalTok{lv4 }\OtherTok{\textless{}{-}}\NormalTok{ df\_comm[df\_comm}\SpecialCharTok{$}\NormalTok{louvain\_comm\_consensus\_r1}\FloatTok{.086} \SpecialCharTok{==} \DecValTok{4}\NormalTok{, ]}\SpecialCharTok{$}\NormalTok{n}
\NormalTok{lv5 }\OtherTok{\textless{}{-}}\NormalTok{ df\_comm[df\_comm}\SpecialCharTok{$}\NormalTok{louvain\_comm\_consensus\_r1}\FloatTok{.086} \SpecialCharTok{==} \DecValTok{5}\NormalTok{, ]}\SpecialCharTok{$}\NormalTok{n}
\NormalTok{lv6 }\OtherTok{\textless{}{-}}\NormalTok{ df\_comm[df\_comm}\SpecialCharTok{$}\NormalTok{louvain\_comm\_consensus\_r1}\FloatTok{.086} \SpecialCharTok{==} \DecValTok{6}\NormalTok{, ]}\SpecialCharTok{$}\NormalTok{n}
\NormalTok{lv7 }\OtherTok{\textless{}{-}}\NormalTok{ df\_comm[df\_comm}\SpecialCharTok{$}\NormalTok{louvain\_comm\_consensus\_r1}\FloatTok{.086} \SpecialCharTok{==} \DecValTok{7}\NormalTok{, ]}\SpecialCharTok{$}\NormalTok{n}
\NormalTok{lv8 }\OtherTok{\textless{}{-}}\NormalTok{ df\_comm[df\_comm}\SpecialCharTok{$}\NormalTok{louvain\_comm\_consensus\_r1}\FloatTok{.086} \SpecialCharTok{==} \DecValTok{8}\NormalTok{, ]}\SpecialCharTok{$}\NormalTok{n}
\NormalTok{lv9 }\OtherTok{\textless{}{-}}\NormalTok{ df\_comm[df\_comm}\SpecialCharTok{$}\NormalTok{louvain\_comm\_consensus\_r1}\FloatTok{.086} \SpecialCharTok{==} \DecValTok{9}\NormalTok{, ]}\SpecialCharTok{$}\NormalTok{n}
\NormalTok{lv10 }\OtherTok{\textless{}{-}}\NormalTok{ df\_comm[df\_comm}\SpecialCharTok{$}\NormalTok{louvain\_comm\_consensus\_r1}\FloatTok{.086} \SpecialCharTok{==} \DecValTok{10}\NormalTok{, ]}\SpecialCharTok{$}\NormalTok{n}
\NormalTok{lv11 }\OtherTok{\textless{}{-}}\NormalTok{ df\_comm[df\_comm}\SpecialCharTok{$}\NormalTok{louvain\_comm\_consensus\_r1}\FloatTok{.086} \SpecialCharTok{==} \DecValTok{11}\NormalTok{, ]}\SpecialCharTok{$}\NormalTok{n}



\NormalTok{color\_nodes\_symp\_df }\OtherTok{\textless{}{-}} \FunctionTok{c}\NormalTok{(}\StringTok{"darkblue"}\NormalTok{, }\StringTok{"limegreen"}\NormalTok{, }\StringTok{"darkorange"}\NormalTok{, }\StringTok{"purple"}\NormalTok{, }\StringTok{"grey"}\NormalTok{, }\StringTok{"yellow"}\NormalTok{, }\StringTok{"black"}\NormalTok{, }\StringTok{"magenta"}\NormalTok{, }\StringTok{"brown"}\NormalTok{, }\StringTok{"darkorchid4"}\NormalTok{ ,}\StringTok{"white"}\NormalTok{ )}


\FunctionTok{qgraph}\NormalTok{(symp\_spear, }\AttributeTok{layout=}\StringTok{"spring"}\NormalTok{, }\AttributeTok{labels =} \FunctionTok{colnames}\NormalTok{(readable\_anx\_dp), }\AttributeTok{sampleSize =} \FunctionTok{nrow}\NormalTok{(df\_anx\_dp\_adj), }\AttributeTok{groups =} \FunctionTok{list}\NormalTok{(}\StringTok{\textquotesingle{}C1 {-} Sui/Appe/Sleep\textquotesingle{}} \OtherTok{=}\NormalTok{ lv1, }\StringTok{\textquotesingle{}C2 {-} Depressed Mood\textquotesingle{}} \OtherTok{=}\NormalTok{ lv2, }\StringTok{\textquotesingle{}C3{-} Conc/Motor/Fatig\textquotesingle{}} \OtherTok{=}\NormalTok{ lv3, }\StringTok{\textquotesingle{}C4 {-} Sleep\textquotesingle{}} \OtherTok{=}\NormalTok{ lv4, }\StringTok{\textquotesingle{}C5 {-} Wort/Dep\textquotesingle{}} \OtherTok{=}\NormalTok{ lv5, }\StringTok{\textquotesingle{}C6 {-} AppeUP\textquotesingle{}} \OtherTok{=}\NormalTok{ lv6, }\StringTok{\textquotesingle{}C7 {-} Dep/Wort\textquotesingle{}} \OtherTok{=}\NormalTok{ lv7, }\StringTok{\textquotesingle{}C8 {-} Dep/Anhe\textquotesingle{}} \OtherTok{=}\NormalTok{ lv8, }\StringTok{\textquotesingle{}C9 {-} Soma/Fear\textquotesingle{}} \OtherTok{=}\NormalTok{ lv9, }\StringTok{\textquotesingle{}C10 {-} Rest\textquotesingle{}} \OtherTok{=}\NormalTok{ lv10, }\StringTok{\textquotesingle{}C11 {-} Anhe/Dep\textquotesingle{}} \OtherTok{=}\NormalTok{ lv11 ), }\AttributeTok{color =}\NormalTok{ color\_nodes\_symp\_df, }\AttributeTok{theme=}\StringTok{"TeamFortress"}\NormalTok{, }\AttributeTok{label.prop=}\FloatTok{0.8}\NormalTok{, }\AttributeTok{vsize=}\DecValTok{2}\NormalTok{, }\AttributeTok{label.cex=}\DecValTok{1}\NormalTok{)}
\end{Highlighting}
\end{Shaded}

\includegraphics{network_analysis_files/figure-latex/unnamed-chunk-32-1.pdf}

\subsection{Creating the Network graph obtained with the best Louvain
community}\label{creating-the-network-graph-obtained-with-the-best-louvain-community}

Each group has the name of the main symptom represented on its
community. For example: the group related to suicide has some items
related to weight loss and change of appetite. This happens likely
because these items tend to cluster together

\begin{Shaded}
\begin{Highlighting}[]
\NormalTok{color\_nodes\_symp\_df }\OtherTok{\textless{}{-}} \FunctionTok{c}\NormalTok{(}\StringTok{"darkblue"}\NormalTok{, }\StringTok{"limegreen"}\NormalTok{, }\StringTok{"darkorange"}\NormalTok{, }\StringTok{"purple"}\NormalTok{, }\StringTok{"grey"}\NormalTok{, }\StringTok{"yellow"}\NormalTok{, }\StringTok{"black"}\NormalTok{, }\StringTok{"magenta"}\NormalTok{, }\StringTok{"brown"}\NormalTok{, }\StringTok{"darkorchid4"}\NormalTok{ ,}\StringTok{"white"}\NormalTok{, }\StringTok{"darkgreen"}\NormalTok{, }\StringTok{"red"}\NormalTok{)}


\NormalTok{anhe }\OtherTok{\textless{}{-}}\NormalTok{ df\_comm[df\_comm}\SpecialCharTok{$}\NormalTok{core\_dsm\_symptoms }\SpecialCharTok{==} \StringTok{\textquotesingle{}anhedonia\textquotesingle{}}\NormalTok{, ]}\SpecialCharTok{$}\NormalTok{n}
\NormalTok{appe }\OtherTok{\textless{}{-}}\NormalTok{ df\_comm[df\_comm}\SpecialCharTok{$}\NormalTok{core\_dsm\_symptoms }\SpecialCharTok{==} \StringTok{\textquotesingle{}apetite/weight\textquotesingle{}}\NormalTok{, ]}\SpecialCharTok{$}\NormalTok{n}
\NormalTok{conc }\OtherTok{\textless{}{-}}\NormalTok{ df\_comm[df\_comm}\SpecialCharTok{$}\NormalTok{core\_dsm\_symptoms }\SpecialCharTok{==} \StringTok{\textquotesingle{}concentration\textquotesingle{}}\NormalTok{, ]}\SpecialCharTok{$}\NormalTok{n}
\NormalTok{dep }\OtherTok{\textless{}{-}}\NormalTok{ df\_comm[df\_comm}\SpecialCharTok{$}\NormalTok{core\_dsm\_symptoms }\SpecialCharTok{==} \StringTok{\textquotesingle{}depressed mood\textquotesingle{}}\NormalTok{, ]}\SpecialCharTok{$}\NormalTok{n}
\NormalTok{fat }\OtherTok{\textless{}{-}}\NormalTok{ df\_comm[df\_comm}\SpecialCharTok{$}\NormalTok{core\_dsm\_symptoms }\SpecialCharTok{==} \StringTok{\textquotesingle{}fatigue\textquotesingle{}}\NormalTok{, ]}\SpecialCharTok{$}\NormalTok{n}
\NormalTok{fear }\OtherTok{\textless{}{-}}\NormalTok{ df\_comm[df\_comm}\SpecialCharTok{$}\NormalTok{core\_dsm\_symptoms }\SpecialCharTok{==} \StringTok{\textquotesingle{}fear/keyed up\textquotesingle{}}\NormalTok{, ]}\SpecialCharTok{$}\NormalTok{n}
\NormalTok{irri }\OtherTok{\textless{}{-}}\NormalTok{ df\_comm[df\_comm}\SpecialCharTok{$}\NormalTok{core\_dsm\_symptoms }\SpecialCharTok{==} \StringTok{\textquotesingle{}Irritability\textquotesingle{}}\NormalTok{, ]}\SpecialCharTok{$}\NormalTok{n}
\NormalTok{motor }\OtherTok{\textless{}{-}}\NormalTok{ df\_comm[df\_comm}\SpecialCharTok{$}\NormalTok{core\_dsm\_symptoms }\SpecialCharTok{==} \StringTok{\textquotesingle{}psychomotor\textquotesingle{}}\NormalTok{, ]}\SpecialCharTok{$}\NormalTok{n}
\NormalTok{rest }\OtherTok{\textless{}{-}}\NormalTok{ df\_comm[df\_comm}\SpecialCharTok{$}\NormalTok{core\_dsm\_symptoms }\SpecialCharTok{==} \StringTok{\textquotesingle{}restlessness/keyed up\textquotesingle{}}\NormalTok{, ]}\SpecialCharTok{$}\NormalTok{n}
\NormalTok{sleep }\OtherTok{\textless{}{-}}\NormalTok{ df\_comm[df\_comm}\SpecialCharTok{$}\NormalTok{core\_dsm\_symptoms }\SpecialCharTok{==} \StringTok{\textquotesingle{}sleep\textquotesingle{}}\NormalTok{, ]}\SpecialCharTok{$}\NormalTok{n}
\NormalTok{soma }\OtherTok{\textless{}{-}}\NormalTok{ df\_comm[df\_comm}\SpecialCharTok{$}\NormalTok{core\_dsm\_symptoms }\SpecialCharTok{==} \StringTok{\textquotesingle{}somatic symptoms\textquotesingle{}}\NormalTok{, ]}\SpecialCharTok{$}\NormalTok{n}
\NormalTok{sui }\OtherTok{\textless{}{-}}\NormalTok{ df\_comm[df\_comm}\SpecialCharTok{$}\NormalTok{core\_dsm\_symptoms }\SpecialCharTok{==} \StringTok{\textquotesingle{}suicide\textquotesingle{}}\NormalTok{, ]}\SpecialCharTok{$}\NormalTok{n}
\NormalTok{wort }\OtherTok{\textless{}{-}}\NormalTok{ df\_comm[df\_comm}\SpecialCharTok{$}\NormalTok{core\_dsm\_symptoms }\SpecialCharTok{==} \StringTok{\textquotesingle{}worthlessness/guilt\textquotesingle{}}\NormalTok{, ]}\SpecialCharTok{$}\NormalTok{n}


\NormalTok{lsymp }\OtherTok{\textless{}{-}} \FunctionTok{list}\NormalTok{(}\StringTok{\textquotesingle{}suicide\textquotesingle{}} \OtherTok{=}\NormalTok{ sui, }\StringTok{\textquotesingle{}depressed\_mood\textquotesingle{}} \OtherTok{=}\NormalTok{ dep, }\StringTok{\textquotesingle{}fatigue\textquotesingle{}} \OtherTok{=}\NormalTok{ fat, }\StringTok{\textquotesingle{}sleep\textquotesingle{}} \OtherTok{=}\NormalTok{ sleep, }\StringTok{\textquotesingle{}worthlessness\textquotesingle{}} \OtherTok{=}\NormalTok{ wort, }\StringTok{\textquotesingle{}appetite\textquotesingle{}} \OtherTok{=}\NormalTok{ appe, }\StringTok{\textquotesingle{}concentration\textquotesingle{}} \OtherTok{=}\NormalTok{ conc , }\StringTok{\textquotesingle{}fear\textquotesingle{}} \OtherTok{=}\NormalTok{ fear, }\StringTok{"irritability"} \OtherTok{=}\NormalTok{ irri , }\StringTok{\textquotesingle{}anhedonia\textquotesingle{}} \OtherTok{=}\NormalTok{ anhe, }\StringTok{\textquotesingle{}psychomotor\textquotesingle{}} \OtherTok{=}\NormalTok{ motor, }\StringTok{\textquotesingle{}restlessness\textquotesingle{}} \OtherTok{=}\NormalTok{ rest, }\StringTok{\textquotesingle{}somatic\textquotesingle{}} \OtherTok{=}\NormalTok{ soma)}


\FunctionTok{qgraph}\NormalTok{(symp\_spear, }\AttributeTok{layout=}\StringTok{"spring"}\NormalTok{, }\AttributeTok{labels =} \FunctionTok{colnames}\NormalTok{(readable\_anx\_dp),}\AttributeTok{sampleSize =} \FunctionTok{nrow}\NormalTok{(df\_anx\_dp\_adj), }\AttributeTok{groups =}\NormalTok{ lsymp, }\AttributeTok{color =}\NormalTok{ color\_nodes\_symp\_df, }\AttributeTok{theme=}\StringTok{"TeamFortress"}\NormalTok{, }\AttributeTok{label.prop=}\FloatTok{0.8}\NormalTok{, }\AttributeTok{vsize=}\DecValTok{2}\NormalTok{, }\AttributeTok{label.cex=}\DecValTok{1}\NormalTok{)}
\end{Highlighting}
\end{Shaded}

\includegraphics{network_analysis_files/figure-latex/unnamed-chunk-33-1.pdf}
\#\# Conclusion

Conectar exploraciones de comunidades con medidas de centralidad. Hacer
inferencias sobre o explorar más en relación a la naturaleza topológica
de síntomas

\end{document}
